\documentclass[12pt,a4paper]{article}
\usepackage[utf8]{inputenc}
\usepackage[turkish]{babel}
\usepackage[T1]{fontenc}
\usepackage{amsmath, amssymb}
\usepackage[hidelinks]{hyperref}
\usepackage{geometry}
\usepackage[most]{tcolorbox}
\usepackage{xcolor}

\geometry{margin=0.8in}

% Turkish babel fix for active characters
\shorthandoff{=}
\shorthandoff{:}

% --- Renk Tanımlamaları ---
\definecolor{cardblue}{RGB}{41, 128, 185}
\definecolor{testblue}{RGB}{30, 144, 255}
\definecolor{ansgreen}{RGB}{34, 139, 34}

% --- Özel Kutular (tcolorbox) ---
\newtcolorbox{learningcard}{
    enhanced,
    colback=white,
    colframe=cardblue!70!white,
    boxrule=0.6pt,
    sharp corners,
    left=15pt, right=15pt, top=12pt, bottom=12pt, middle=15pt,
    before skip=2em, after skip=2em,
    segmentation style={cardblue!10!white, line width=0.5pt},
    shadow={2pt}{-2pt}{0pt}{black!5!white}
}

\newtcolorbox{testcard}[1]{
    enhanced,
    colback=white,
    colframe=testblue!70!white,
    boxrule=0.6pt,
    arc=0pt,
    outer arc=0pt,
    title={\textbf{#1}},
    coltitle=white,
    colbacktitle=testblue!80!black,
    fonttitle=\large\bfseries,
    left=10pt, right=10pt, top=10pt, bottom=10pt,
    segmentation style={draw=testblue!30!white, solid, linewidth=1pt},
    lower separated=true
}

% --- Sayaçlar ve Başlıklar ---
\newcounter{soruno}
\newcommand{\SoruHeader}{\stepcounter{soruno}{\textbf{\color{red!80!black}Soru \thesoruno:}}\ }
\newcommand{\CevapHeader}{{\textbf{\color{green!60!black}Cevap:}}\ }

% --- Kapak Bilgileri ---
\title{\Huge \textbf{Makine Öğrenmesine Giriş}}
\author{Oğuz Taşdemir}
\date{\today}

\begin{document}

\maketitle
\newpage

\tableofcontents
\newpage

% ==============================================================================
% BÖLÜM 1: EL KİTABI
% ==============================================================================
\section*{Giriş}
\addcontentsline{toc}{section}{Giriş}
Bu çalışma, İstanbul Medeniyet Üniversitesi bünyesinde verilen Makine Öğrenmesine Giriş dersi kapsamında hazırlanan kapsamlı bir ders rehberidir. Temel amaç, dönem boyunca işlenen teorik konuları, algoritmaların çalışma mantığını ve uygulama stratejilerini bir araya getirerek öğrenciye bütünsel bir bakış açısı sunmaktır. 

Döküman; veri hazırlığından derin öğrenmeye kadar uzanan geniş bir yelpazeyi, hem akademik disiplinle hem de anlaşılır bir dille ele almaktadır. Rehberin, sınav hazırlık sürecinde ve projelerde temel bir başvuru kaynağı olması hedeflenmiştir.

\newpage
\section{Makine Öğrenmesinin Temelleri}
Makine öğrenmesi algoritmaları, öğrenme biçimlerine göre üç ana kategoriye ayrılır:

\begin{itemize}
    \item \textbf{Denetimli (Supervised) Öğrenme:} Veride hem girdi hem de hedef etiket (label) bulunur. Model, girdiler ile çıktılar arasındaki ilişkiyi öğrenir (Örn: Sınıflandırma, Regresyon).
    \item \textbf{Denetimsiz (Unsupervised) Öğrenme:} Veride etiket yoktur. Algoritma, verinin kendi içindeki gizli yapıları veya benzerlikleri keşfeder (Örn: Kümeleme, Boyut Düşürme).
    \item \textbf{Pekiştirmeli (Reinforcement) Öğrenme:} Bir ajanın (agent), bir ortamda (environment) en yüksek ödülü almak için deneme-yanılma yoluyla öğrenmesidir.
\end{itemize}

\subsection{Makine Öğrenmesi Nedir?}
Bilgisayarların açıkça programlanmadan veriden öğrenmesini sağlayan bir yapay zeka alt dalıdır. Klasik programlamada girdi ve kurallar verilip çıktı alınırken; makine öğrenmesinde girdi ve çıktılar verilerek bu ikisi arasındaki ilişkiyi temsil eden model (kurallar seti) öğrenilir.

Matematiksel olarak bir öğrenme süreci, bir performans ölçütünü (hata payı gibi) optimize edecek şekilde parametrelerin ($\theta$) güncellenmesidir:
\begin{equation}
    y = g(x | \theta)
\end{equation}
Burada $g$ öğrenilen fonksiyonu, $x$ girdiyi, $\theta$ ise modelin iç parametrelerini temsil eder. Amaç, gerçek değer ile tahmin arasındaki farkı (kayıp/loss) minimize etmektir. Optimizasyon sürecinde \textbf{Gradyan İnişi (Gradient Descent)} kullanılır. Burada öğrenme hızı ($\alpha$) adımların büyüklüğünü belirler; $\alpha$ çok büyükse model \textbf{Iraksama} (divergence) yapabilir, çok küçükse \textbf{Yerel Minimum} (local minimum) noktasına takılıp kalabilir veya \textbf{Yakınsama} (convergence) süreci çok uzayabilir.

\subsubsection{Modern Optimizerlar: Adam}
\textbf{Adam (Adaptive Moment Estimation):} Momentum (geçmiş gradyanların ortalaması) ve RMSProp (geçmiş gradyanların karelerinin hareketli ortalaması) tekniklerini birleştirir. Her parametre için farklı bir öğrenme hızı belirleyerek eğitimi hızlandırır ve daha kararlı hale getirir.

\subsubsection{Denetimli, Denetimsiz ve Yarı-Denetimli Öğrenme}
\begin{itemize}
    \item \textbf{Denetimli Öğrenme (Supervised Learning):} Veri setinde her girdinin bir karşılık gelen etiketi (label) bulunur. İki ana alt başlığa ayrılır:
    \begin{itemize}
        \item \textbf{Sınıflandırma (Classification):} Çıktı kategoriktir (Örn: Spam/Değil, Kredi Riski: Yüksek/Düşük).
        \item \textbf{Regresyon (Regression):} Çıktı süreklidir (Örn: Ev fiyatı tahmini, hisse senedi değeri).
    \end{itemize}
    \item \textbf{Denetimsiz Öğrenme (Unsupervised Learning):} Etiketsiz verilerdeki gizli yapıları ve örüntüleri bulmayı hedefler. Temel uygulamaları \textbf{Kümeleme (Clustering)} ve \textbf{Boyut Düşürme}'dir.
    \item \textbf{Yarı-Denetimli Öğrenme:} Az miktarda etiketli veri ile çok büyük miktarda etiketsiz verinin birlikte kullanıldığı hibrit yöntemdir.
\end{itemize}

\subsubsection{Takviyeli Öğrenme (Reinforcement Learning) Temelleri}
Bir ajanın (agent), bir ortamda (environment) eylemler gerçekleştirerek en yüksek ödülü (reward) toplamaya çalıştığı öğrenme türüdür. "Gecikmiş Ödül" kavramı kritiktir; ajan anlık kazanç yerine uzun vadeli stratejiye odaklanır (Örn: Satranç motorları, otonom sürüş).

\subsubsection{Temel Teoremler: No Free Lunch ve Inductive Bias}
\begin{itemize}
    \item \textbf{No Free Lunch (Bedava Öğle Yemeği Yok):} Matematiksel olarak, tüm olası problemler üzerinde diğerlerinden her zaman daha iyi performans gösteren tek bir "en iyi" algoritma yoktur. Model başarısı tamamen veri setinin yapısına ve probleme uygunluğuna bağlıdır.
    \item \textbf{Endüktif Önyargı (Inductive Bias):} Bir modelin eğitim verisinin ötesine geçip genelleme yapabilmesi için sahip olduğu varsayımlar bütünüdür. Örneğin, doğrusal regresyonun verideki ilişkinin doğrusal olduğunu varsayması bir önyargıdır.
\end{itemize}

\subsubsection{Makine Öğrenmesi İş Akışı (Pipeline)}
Sıfırdan başlayan bir projede takip edilen standart adımlar şöyledir:
\begin{enumerate}
    \item \textbf{Veri Toplama:} Ham verinin (CSV, SQL, API vb.) sisteme alınması.
    \item \textbf{Veri Ön İşleme:} Eksik değerlerin doldurulması ve ölçekleme (Bölüm 3).
    \item \textbf{Model Seçimi ve Eğitim:} Probleme uygun algoritmanın seçilip verilerle "eğitilmesi".
    \item \textbf{Değerlendirme:} Modelin hiç görmediği "test verisi" ile başarısının ölçülmesi.
    \item \textbf{Optimizasyon:} Hiperparametreleri değiştirerek en iyi sonucun aranması.
\end{enumerate}

\newpage
\subsection{Hata ve Performans Analizi}
Bir modelin başarısı, toplam hatanın bileşenlerine ayrıştırılmasıyla analiz edilir. Toplam hata (MSE), \textbf{İndirgenebilir Hata} ve \textbf{İndirgenemez Hata} ($\sigma^2$) olarak ikiye ayrılır.

\subsubsection{Bias-Variance Ödünleşimi (Trade-off)}
\begin{itemize}
    \item \textbf{Bias (Yanlılık):} Modelin verideki temel ilişkiyi kavrayamayacak kadar basit olması durumudur (Underfitting).
    \item \textbf{Variance (Varyans):} Modelin eğitim verisindeki gürültüyü ezberleyerek aşırı esnek hale gelmesi durumudur (Overfitting).
\end{itemize}
İdeal bir çalışma, her iki hatanın da dengelendiği noktayı bulmayı hedefler.

\subsubsection{Hata Bileşenleri: $MSE = Bias^2(f) + Var(f) + \sigma^2$}
Hata formülündeki $\sigma^2$, veri setindeki doğal gürültüdür ve hiçbir model tarafından yok edilemez. Mühendislik çabası, $Bias^2 + Var$ toplamını minimize etmeye yöneliktir.

\subsubsection{Sınıflandırma Hata Oranı ($E$)}
Sınıflandırma modellerinde temel başarı ölçütü olan hata oranı, tahmin edilen etiketin ($\hat{v}_i$) gerçek etiketten ($v_i$) farklı olduğu durumların toplam örneğe oranıdır:
\begin{equation}
    E = \frac{1}{n} \sum_{i=1}^{n} L(\hat{v}_i, v_i) \quad \text{burada } L=1 \text{ eğer } \hat{v}_i \neq v_i, \text{ aksi halde } 0
\end{equation}
Bu formül, dökümanın ilerleyen bölümlerinde detaylandırılacak olan "Confusion Matrix" (Karmaşıklık Matrisi) üzerindeki köşegen dışı elemanların toplamını temsil eder.

\subsubsection{Aşırı Öğrenme (Overfitting) ve Eksik Öğrenme (Underfitting)}
\begin{itemize}
    \item \textbf{Overfitting (Aşırı Uyum):} Eğitim hatası çok düşük, ancak test hatası çok yüksektir. Model verideki gürültüyü bile ezberlemiştir.
    \\ \textbf{Çözüm Stratejileri:} 
    \begin{itemize}
        \item Daha fazla veri toplamak.
        \item \textbf{Düzenlileştirme (Regularization):} Modelun katsayılarına ceza puanı vererek (L1/L2 gibi yöntemlerle) aşırı büyük değerler almasını engellemek.
        \item \textbf{Dropout:} Eğitim sırasında nöronların bir kısmını rastgele devre dışı bırakarak modelin belirli yollara bağımlı kalmasını önlemek.
        \item Model karmaşıklığını azaltmak (katman veya özellik sayısını düşürmek).
    \end{itemize}
    
    \item \textbf{Underfitting (Yetersiz Uyum):} Hem eğitim hem test hatası yüksektir. Model veriye dair temel yapıyı bile kavrayamamıştır.
    \\ \textbf{Çözüm Stratejileri:} 
    \begin{itemize}
        \item Daha karmaşık/derin modeller kullanmak.
        \item \textbf{Özellik Mühendisliği (Feature Engineering):} Yeni ve daha açıklayıcı özellikler türetmek.
        \item Düzenlileştirme (ceza) miktarını azaltmak.
    \end{itemize}
\end{itemize}

\subsection{Veri Yönetimi Stratejileri}
\subsubsection{Eğitim, Test ve Doğrulama (Validation) Setleri}
Modelin genelleme yeteneğini ölçmek için veri genellikle \%80 Eğitim, \%20 Test olarak bölünür. Hiperparametre optimizasyonu için ise Eğitim setinden bir "Doğrulama (Validation)" seti ayrılır.

\subsubsection{Model Seçimi: AIC ve BIC Kriterleri}
Hangi modelin daha iyi olduğunu belirlemek için hata payı ile parametre sayısını dengeleyen kriterlerdir:
\begin{itemize}
    \item \textbf{AIC (Akaike) ve BIC (Bayesian):} Değer ne kadar \textbf{küçükse} model o kadar iyidir. Bu kriterler, modele eklenen her gereksiz parametre için modeli cezalandırır.
\end{itemize}
\begin{itemize}
    \item \textbf{Parametre:} Modelin eğitim sırasında veriden kendi öğrendiği değerlerdir (Örn: Regresyondaki katsayılar/ağırlıklar).
    \item \textbf{Hiperparametre:} Eğitim başlamadan önce \textbf{bizim} dışarıdan belirlediğimiz ayar düğmeleridir (Örn: k-NN'deki $k$ sayısı, Sinir Ağlarındaki katman sayısı).
\end{itemize}

\subsubsection{Veri Sızıntısı (Data Leakage) Kavramı}
Test setindeki bilginin model eğitim aşamasında yanlışlıkla kullanılmasıdır. \textbf{Hayati Uyarı:} Bu durum, modelin gerçekte olduğundan çok daha başarılı görünmesine (sahte başarı) yol açar ve model gerçek veriye çıktığında çöker.

\newpage
\section{Veri Hazırlığı ve Özellik Mühendisliği}

\subsection{Veri Ön İşleme (Preprocessing)}
Ham verinin model tarafından işlenebilir "Sadık Embedding (Vektör)" yapılarına dönüştürülme sürecidir. Bir modelin başarısı, veriyi nasıl "gördüğüne" doğrudan bağlıdır.

\subsubsection{Ölçekleme: Standartlaştırma (Z-Score) ve Normalizasyon}
\begin{itemize}
    \item \textbf{Z-Skoru Standartlaştırma:} Verinin ortalamasını 0, standart sapmasını 1 yapar. Özelliklerin farklı ölçeklerde (Örn: Yaş ve Gelir) olmasından kaynaklanan "baskınlık" sorununu çözer.
    \item \textbf{Min-Max Normalizasyon:} Veriyi $[0, 1]$ aralığına sıkıştırır. Görüntü işleme gibi belirli sınırlar gerektiren alanlarda tercih edilir.
\end{itemize}

\subsubsection{Eksik Değer Analizi ve İmputasyon}
Veri setindeki boşluklar (NaN) model eğitimini durdurabilir.
\begin{itemize}
    \item \textbf{Temel Müdahaleler:} Satır silme (veri kaybı azsa) veya basit atama (ortalama/mod kullanımı).
    \item \textbf{Akıllı Tahminleme:} Eksik değerin, diğer değişkenler kullanılarak bir regresyon modeli veya PCA (Temel Bileşen Analizi) ile tahmin edilmesidir.
\end{itemize}

\subsubsection{Dağılım Yönetimi ve Dönüşümler}
\begin{itemize}
    \item \textbf{Log ve Gama Dönüşümü:} Web trafiği gibi çok geniş ölçekte değişen "çarpık" (skewed) verileri normale yaklaştırmak için kullanılır.
    \item \textbf{Winsorizing (Kırpma):} Aykırı değerlerin (outliers) etkisini azaltmak için veriyi belirli yüzdelik dilimler (Örn: \%5 ve \%95) arasında sabitleme işlemidir.
    \item \textbf{One-Hot Encoding:} Kategorik verileri (Elma, Armut, Muz) aralarında büyüklük ilişkisi kurmadan $[1, 0, 0]$ gibi birim vektörlere dönüştürür.
\end{itemize}

\subsection{Özellik Seçimi (Feature Selection)}
Gereksiz özelliklerin elenmesi; aşırı uyumu azaltır, hesap verimliliğini ve yorumlanabilirliği artırır.

\begin{enumerate}
    \item \textbf{Filtre Yöntemleri (Filter):} Modelden bağımsız, istatistiksel skorlara dayanır. Hızlıdır. (Örn: Pearson Korelasyonu, Chi-Kare, Mutual Information).
    \begin{itemize}
        \item \textbf{ANOVA F-Testi:} Sınıf ortalamaları arasındaki farkın, sınıf içi varyansa oranını ($F$ değeri) ölçer. Yüksek $F$ değeri, özelliğin ayırt edici gücünün yüksek olduğunu gösterir.
    \end{itemize}
    \item \textbf{Sarma Yöntemleri (Wrapper):} Özellik alt kümelerini bir model üzerinde test eder. Yavaştır ama doğrudur. (Örn: Forward Selection, RFE - Recursive Feature Elimination).
    \item \textbf{Gömülü Yöntemler (Embedded):} Özellik seçimi eğitimle birlikte gerçekleşir. (Örn: LASSO - L1 Düzenlileştirme, Karar Ağacı Önemi).
\end{enumerate}

\subsection{Boyutların Laneti (Curse of Dimensionality)}
Özellik sayısı arttıkça, veriyi temsil etmek için gereken örneklem miktarı üstel olarak artar. Bu durum geometrik bir seyrekliğe yol açar; yüksek boyutlarda veri noktaları birbirine çok uzak kalır ve benzerlik ölçümleri (mesafeler) anlamını yitirmeye başlar. Bu nedenle, modelin genelleme yeteneğini korumak için boyut düşürme veya özellik seçimi kritik önem taşır.

\subsection{Mesafe ve Benzerlik Ölçütleri}
Bir ölçümün matematiksel olarak \textbf{metrik} sayılabilmesi için üç temel kuralı sağlaması gerekir: Negatif olmama ($d(x,y) \geq 0$), simetri ($d(x,y)=d(y,x)$) ve \textbf{Üçgen Eşitsizliği} ($d(x,y) \leq d(x,z) + d(z,y)$).

\begin{itemize}
    \item \textbf{L1 (Manhattan) Mesafesi:} Izgara tabanlı şehir yapısındaki hareket gibi mutlak farkların toplamıdır. Yüksek boyutlu uzaylarda ve aykırı değerlere karşı Öklid'e göre daha dirençlidir.
    \item \textbf{L2 (Öklid) Mesafesi:} İki nokta arasındaki en kısa doğrusal mesafedir. Standart yaklaşımdır.
    \item \textbf{Kosinüs Benzerliği:} Vektörlerin büyüklüğüne değil, \textbf{yönüne ve aralarındaki açıya} odaklanır. Metin madenciliğinde (NLP) yaygındır.
    \item \textbf{Jaccard İndeksi:} Kümelerin ne kadar örtüştüğünü ölçer. İkili (binary) veriler için idealdir.
\end{itemize}

\newpage
\section{Regresyon Modelleri}
Regresyon; bağımsız değişkenler ($X$) ile sürekli bir bağımlı değişken ($y$) arasındaki ilişkiyi öğrenen denetimli bir öğrenme yöntemidir. Temel amaç, verilen bir girdi için sayısal bir tahmin üretmektir.

\subsection{Doğrusal ve Çoklu Regresyon}
\subsubsection{En Küçük Kareler (OLS) ve Normal Denklem}
OLS yöntemi, tahmin edilen değerler ile gerçek değerler arasındaki artıkların (residuals) kareleri toplamını (SSE) minimize etmeyi hedefler. Analitik çözüm olan Normal Denklem şu şekildedir:
\begin{equation}
    \beta = (X^T X)^{-1} X^T y
\end{equation}
\textbf{Güven Notu:} Bu formül işin matematiksel temelidir. Uygulamada bu matris işlemlerini Scikit-Learn gibi kütüphaneler otomatik olarak, milisaniyeler içinde yapar.
OLS, veri temiz ve az boyutlu olduğunda hızlıdır ancak aykırı değerlere (outliers) karşı oldukça duyarlıdır. Ayrıca bağımsız değişkenler arasında yüksek korelasyon olması (\textbf{Çoklu Doğrusallık / Multicollinearity}) katsayı tahminlerinin varyansını artırarak modeli kararsızlaştırır. Analitik çözümün geçerli olması için özelliklerin birbirinden bağımsız olması idealdir.

\subsubsection{Polinomsal Regresyon}
Doğrusal olmayan (eğrisel) ilişkileri modellemek için özellikleri yüksek dereceli terimlere ($x^2, x^3$) dönüştürerek doğrusal modellerle eğitme işlemidir. Aşırı karmaşık dereceler varyansı artırır.

\subsubsection{RANSAC: Sağlam (Robust) Regresyon}
Verideki aykırı değerlerin modeli bozmasını engellemek için kullanılır. Algoritma, veriden rastgele küçük bir alt küme (inliers) seçerek model kurar ve en çok veriyi açıklayan temiz modeli belirler. Bu, özellikle veri setinde çok fazla hatalı ölçüm veya uç değer varsa standart OLS'den çok daha güvenilir sonuçlar üretir. \textbf{Aykırı değerlere karşı dayanıklı (Robust)} bir yaklaşımdır.

\subsection{Düzenlileştirme (Regularization) ve Seyreklik}
Modelin katsayılarının aşırı büyümesini engelleyerek aşırı uyumu (overfitting) önleyen yöntemlerdir.

\begin{itemize}
    \item \textbf{Ridge Regresyon (L2):} Hata fonksiyonuna katsayıların karelerinin toplamını ekler ($\alpha \cdot \sum \beta_j^2$). Katsayıları sıfıra yaklaştırır ama tam sıfır yapmaz. Çoklu doğrusallık (multicollinearity) durumunda etkilidir.
    \item \textbf{Lasso Regresyon (L1):} Hata fonksiyonuna katsayıların mutlak değerlerinin toplamını ekler ($\alpha \cdot \sum |\beta_j|$). Bazı katsayıları tam olarak sıfıra yaparak \textbf{Özellik Seçimi} sağlar. Seyrek (sparse) modeller üretir.
    \item \textbf{Elastic Net:} L1 ve L2 düzenlileştirmelerini birleştiren hibrit yaklaşımdır. Gruplu değişkenlerin olduğu senaryolarda Lasso'dan daha iyi performans gösterir.
\end{itemize}

\subsection{Hata Ölçütleri ve Yorumlama}
\begin{itemize}
    \item \textbf{MAE (Mean Absolute Error):} Hataların mutlak ortalamasıdır. Aykırı değerlere karşı stabildir ve doğrusal bir ceza uygular.
    \item \textbf{MSE (Mean Squared Error):} Hataların karesinin ortalamasıdır. Büyük hataları karesel olarak büyüttüğü için çok daha sert cezalandırır.
    \item \textbf{RMSE (Root MSE):} MSE'nin kareköküdür. Hata birimini orijinal hedef birimine döndürerek yorumlamayı kolaylaştırır.
    \item \textbf{MAPE (Mean Absolute Percentage Error):} Hata payını yüzde olarak verir. \textbf{Kritik Zayıflık:} Gerçek değerlerin 0 olduğu durumlarda tanımsız hale gelir; ayrıca düşük değerlerde çok yüksek yüzde hataları üretebilir.
\end{itemize}

\subsubsection{R-Kare ve Düzeltilmiş R-Kare}
\textbf{R-Kare ($R^2$):} Bağımsız değişkenlerin, hedefteki varyansın yüzde kaçını açıkladığını gösterir.
\\ \textbf{Düzeltilmiş R-Kare ($Adjusted R^2$):} Modele eklenen gereksiz değişkenler için ceza uygular. Standart $R^2$ asla düşmezken, $Adjusted R^2$ model kalitesi artmadıkça yükselmez; böylece sahte başarının önüne geçilir.

\newpage
\section{Sınıflandırma Modelleri}
Sınıflandırma, veri noktalarını önceden tanımlanmış kategorilere (etiketlere) atama sürecidir. Çıktı değişkeni süreksiz ve kategoriktir.

\subsection{Lojistik Tahmin}
Lojistik regresyon, ismine rağmen bir sınıflandırma algoritmasıdır. Çıktıyı $[0, 1]$ arasına sıkıştırmak için \textbf{Sigmoid (Lojistik)} fonksiyonunu kullanır. İkili (binary) sınıflandırma problemlerinde (Örn: Hasta/Sağlıklı) temel yöntemdir.

\subsection{Örnek Tabanlı Öğrenme: k-NN}
k-En Yakın Komşu (k-NN) algoritması, \textbf{Tembel Öğrenme (Lazy Learning)} olarak bilinir. Eğitim aşamasında model kurmaz, tüm hesaplamayı tahmin anında yapar. Bir noktanın sınıfını, ona en yakın $k$ adet komşusunun oylarıyla belirler.
\\ \textbf{Kritik Detay:} $k$ sayısının çok küçük seçilmesi aşırı uyuma (overfitting), çok büyük seçilmesi ise yetersiz uyuma (underfitting) yol açar. Aramayı hızlandırmak için \textbf{KD-Tree} (k-boyutlu ağaç) yapısı kullanılır; bu yapı uzayı eksenlere dik bölerek en yakın komşu arama maliyetini düşürür. Mesafe tabanlı olduğu için \textbf{Scaling} (ölçekleme) yapılması zorunludur.

\subsection{Destek Vektör Makineleri (SVM)}
SVM, sınıflar arasındaki boşluğu (\textbf{Margin}) maksimize eden en uygun karar hiperdüzlemini bulmayı hedefler.
\begin{itemize}
    \item \textbf{Kernel Trick (Çekirdek Yöntemi):} Doğrusal olarak ayrılamayan verileri yüksek boyutlu bir uzaya taşıyarak ayrılabilir hale getirir.
    \item \textbf{Destek Vektörleri:} Karar sınırını belirleyen, sınırlara en yakın olan veri noktalarıdır.
\end{itemize}

Bayes Teoremi'ne dayanır. Özelliklerin sınıf verildiğinde birbirinden tamamen \textbf{bağımsız} olduğunu varsaydığı için "Naive" (Saf) olarak adlandırılır.
\\ \textbf{Türleri:}
\begin{itemize}
    \item \textbf{Gaussian Naive Bayes:} Özelliklerin normal (Gauss) dağılıma sahip olduğu sürekli veriler için kullanılır.
    \item \textbf{Bernoulli Naive Bayes:} Özelliklerin ikili (binary) olduğu durumlar için (Örn: Kelime var/yok).
    \item \textbf{Multinomial Naive Bayes:} Kelime frekansları (sayımları) gibi ayrık veriler için idealdir.
\end{itemize}
\textbf{Laplace Smoothing:} Veride olmayan bir özelliğin olasılığının 0 olmasını engellemek için tüm sayımlara 1 eklenmesidir.

\subsection{Karar Yapıları ve Topluluklar (Ensembles)}
\begin{itemize}
    \item \textbf{Karar Ağaçları:} Veriyi \textbf{Entropi} (logaritmik) veya \textbf{Gini Safsızlığı} (daha hızlı, olasılıksal) ölçütlerine göre dallara ayırır. Gini genellikle büyük verilerde hız avantajı sağlar.
    \item \textbf{Random Forest (Bagging):} Çok sayıda karar ağacını hem \textbf{Bootstrap} (rastgele örneklem) hem de \textbf{Feature Randomness} (her dalda rastgele özellik seçimi) ile eğiterek varyansı düşürür.
    \item \textbf{Gradient Boosting (Boosting):} Ağaçları ardışıl (sequential) olarak eğitir; her yeni ağaç bir önceki ağacın hatalarını (residuals) minimize etmeye odaklanır.
\end{itemize}

\newpage
\section{Kümeleme ve Yapı Bulma}
Kümeleme, veri setindeki benzer örnekleri gruplandırmayı amaçlayan denetimsiz bir öğrenme yöntemidir. Önceden tanımlanmış bir etiket (label) yoktur; yapı verinin kendi özelliklerinden keşfedilir.

\subsection{Merkez Tabanlı Kümeleme: K-Means}
Veriyi $K$ adet kümeye ayıran iteratif bir algoritmadır. Her nokta en yakın küme merkezine (centroid) atanır ve merkezler her adımda güncellenir.
\begin{itemize}
    \item \textbf{Inertia:} Veri noktalarının kendi küme merkezlerine olan uzaklıklarının kareleri toplamıdır. Düşük inertia, daha sıkı kümeler anlamına gelir.
    \item \textbf{Dirsek (Elbow) Yöntemi:} Optimal küme sayısını ($K$) belirlemek için kullanılır. Inertia'daki düşüşün "dirsek" yaptığı nokta seçilir.
    \item \textbf{Silhouette Skoru:} Bir noktanın kendi kümesine ne kadar benzer, komşu kümelere ne kadar uzak olduğunu ölçer. $+1$ mükemmel kümelemeyi, $-1$ hatalı atamayı gösterir.
\end{itemize}

\subsection{Hiyerarşik ve Yoğunluk Tabanlı Yapılar}
\begin{itemize}
    \item \textbf{Hiyerarşik Kümeleme:} Veriyi ağaç benzeri bir yapıda (\textbf{Dendrogram}) birleştirir. \textbf{Single Linkage} (en yakın komşu) dairesel olmayan yapıları yakalayabilir ancak 'zincirleme' (chaining) etkisine yatkındır. \textbf{Complete Linkage} (en uzak komşu) daha kompakt kümeler üretir.
    \item \textbf{DBSCAN:} Yoğunluk tabanlıdır. \textbf{Eps} (yarıçap) ve \textbf{MinPts} (min. komşu) parametrelerini kullanır. Karmaşık şekilli kümeleri bulabilir ve gürültüyü eleyebilir.
\end{itemize}
\textbf{Sanity Check:} Kümeleme sonuçlarının iş mantığına (Örn: Pazar segmentasyonu) uygunluğunun uzman kontrolüyle doğrulanmasıdır.

\textbf{Gauss Karışım Modelleri (GMM):} Verinin farklı olasılık dağılımlarının birleşimi olduğunu varsayar ve "yumuşak" (\textbf{Soft Clustering} - olasılıksal) kümeleme yapar. K-Means gibi yöntemler "sert" (\textbf{Hard Clustering}) kümeleme yapar. GMM, \textbf{EM (Expectation-Maximization)} algoritması ile optimize edilir.

\newpage
\section{Boyut Düşürme ve Veri Sıkıştırma}
Yüksek boyutlu verilerdeki gereksiz varyansı temizleyerek veriyi daha düşük boyutlu bir alt uzaya izdüşürme işlemidir.

\subsection{Doğrusal İzdüşüm Yöntemleri}
\begin{itemize}
    \item \textbf{PCA (Temel Bileşen Analizi):} Verideki maksimum varyansı koruyan ortogonal bileşenler üretir. \textbf{Açıklanan Varyans Oranı (EVR)}, her bileşenin toplam bilginin ne kadarını taşıdığını yüzdece ifade eder.
    \item \textbf{LDA (Doğrusal Diskriminant Analizi):} PCA'nın aksine denetimli (supervised) bir yöntemdir. Sınıf ayrımını (separability) maksimize etmeyi hedefler. Sınıf sayısı $C$ ise, en fazla $C-1$ adet yeni bileşen üretebilir.
\end{itemize}

\subsection{Karmaşık ve Doğrusal Olmayan Dönüşümler}
\begin{itemize}
    \item \textbf{Otoenkodörler (Autoencoders):} Veriyi \textbf{Bottleneck} (Darboğaz) katmanında sıkıştırıp tekrar açan sinir ağlarıdır. Burada oluşan düşük boyutlu temsile \textbf{Latent Space} (Gizli Uzay) adı verilir. Gürültü temizleme (\textbf{Denoising}) için de kullanılır.
    \item \textbf{t-SNE ve UMAP:} Verinin 2D/3D görselleştirilmesi için kullanılan, komşuluk ilişkilerini koruyan doğrusal olmayan algoritmalardır.
\end{itemize}

\newpage
\section{Yapay Sinir Ağları ve Derin Öğrenme}
İnsan beynindeki nöron yapısından esinlenen, karmaşık doğrusal olmayan ilişkileri öğrenebilen katmanlı yapılardır.

\subsection{Nöron Yapısı ve Aktivasyon Fonksiyonları}
Bir yapay nöron; girdileri (\textbf{Weights}), bir yanlılık değerini (\textbf{Bias}) alır ve bunları bir aktivasyon fonksiyonundan geçirerek çıktı üretir.
\begin{itemize}
    \item \textbf{ReLU:} Negatif değerleri sıfırlar, doğrusal olmayan yapıyı hızlı şekilde öğrenir.
    \item \textbf{Sigmoid/Softmax:} Çıktıyı olasılık değerine dönüştürür.
\end{itemize}

\subsection{Eğitim Süreci: Geri Yayılım (Backpropagation)}
Çıkış katmanında oluşan hatanın, \textbf{Zincir Kuralı (Chain Rule)} kullanılarak geriye doğru ağa dağıtılması ve ağırlıkların güncellenmesidir. Bu süreçte her parametre, hataya katkısı oranında değiştirilir. \textbf{XOR Problemi} gibi doğrusal olarak ayrılamayan verileri çözebilmek için en az bir gizli katman (\textbf{Hidden Layer}) kullanılır. \textbf{Evrensel Yaklaşım Teoremi}'ne göre, yeterince geniş bir gizli katman her türlü sürekli fonksiyonu modelleyebilir.

\subsection{Optimizasyon ve Düzenlileştirme}
Derin ağlarda \textbf{Gradyan Kaybolması} (Vanishing Gradient) sorununu aşmak için \textbf{Batch Normalization} (katman girişlerini normalize ederek \textbf{Internal Covariate Shift} etkisini azaltır) gibi teknikler kullanılır. Ayrıca aşırı öğrenmeyi önlemek için bazı nöronların rastgele kapatıldığı \textbf{Dropout} yöntemi yaygındır.

\newpage
\section{Model Değerlendirme ve Optimizasyon}

\subsection{Doğrulama ve Çapraz Doğrulama (Cross-Validation)}
\begin{itemize}
    \item \textbf{k-Fold CV:} Veriyi $k$ parçaya böler. Genellikle $k=5$ veya $k=10$ seçilir.
    \item \textbf{Bootstrap (Önyükleme):} Veriden \textbf{yerine koyarak} (with replacement) örnek çekilmesidir.
    \item \textbf{LOOCV (Leave-One-Out):} Küçük verilerde hassastır ancak yüksek hesaplama maliyeti dezavantajıdır.
    \item \textbf{Nested (İç İçe) CV:} Aynı anda hem hiperparametre ayarlamak hem de tarafsız performans ölçmek için kullanılır.
\end{itemize}

\subsection{Hiperparametre Optimizasyonu Stratejileri}
\begin{itemize}
    \item \textbf{Grid Search (Izgara Araması):} Belirlenen tüm kombinasyonları dener. Hesaplama maliyeti $O(L^F)$ formülüyle ifade edilir ($L$: Seviye, $F$: Faktör).
    \item \textbf{Random Search (Rastgele Arama):} Yüksek boyutlu uzaylarda, önemli olan parametreleri daha verimli taradığı için Grid Search'ten daha başarılı olabilir.
    \item \textbf{Simulated Annealing:} "Soğuma" dinamiğiyle yerel minimumlardan kaçmayı hedefler.
    \item \textbf{Genetik Algoritmalar:} \textbf{Mutasyon ve Çaprazlama} mekanizmalarıyla en iyi parametreleri arar.
\end{itemize}

\subsection{Başarı Metrikleri}
\begin{itemize}
    \item \textbf{Karmaşıklık Matrisi:} \textbf{TP} (Doğru Pozitif), \textbf{TN} (Doğru Negatif), \textbf{FP} (Tip I Hata - Yalancı Alarm) ve \textbf{FN} (Tip II Hata - Kaçırılan Vaka) değerlerini gösterir.
    \item \textbf{F1 Skoru:} Precision ve Recall'un \textbf{Harmonik Ortalamasıdır}. Harmonik ortalama, aritmetik ortalamadan farklı olarak düşük olan değere karşı çok daha hassastır; yani kesinlik veya duyarlılıktan biri çok düşükse F1 skoru çöker.
    \item \textbf{ROC-AUC:} Modelin sınıfları ayırma kabiliyetini [0, 1] aralığında ölçer. $AUC=0.5$ rastgele tahmini, $AUC=1$ mükemmel ayrımı ifade eder.
    \item \textbf{Neyman-Pearson Kriteri:} Hata maliyetlerinin farklı olduğu durumlarda kullanılır. Toplam hata $E = \lambda \frac{C_{fn}}{n} + \frac{C_{fp}}{n}$ formülüyle hesaplanır. Burada $\lambda > 1$ seçilmesi, yanlış negatiflerin (vaka kaçırmanın) daha maliyetli kabul edildiği anlamına gelir.
\end{itemize}

\newpage
\section{Özet Karşılaştırmalar ve Kritik İpuçları}
Bu bölüm, öğrenme kartlarında ve sınav bankasında sıkça karşınıza çıkabilecek "ayırt edici" detayları özetlemektedir:

\subsection{Regresyon ve Düzenlileştirme}
\begin{itemize}
    \item \textbf{Ridge (L2) vs. Lasso (L1):} Lasso katsayıları tam olarak sıfıra çekebilir ve otomatik özellik seçimi yapar. Ridge ise katsayıları küçültür ama sıfırlamaz.
    \item \textbf{R-Kare (R2):} Modelin veriyi açıklama oranını gösterir. Ayarlanmış R2 ise modele eklenen gereksiz özelliklerin başarısını cezalandırır.
\end{itemize}

\subsection{Mesafe ve Sınıflandırma}
\begin{itemize}
    \item \textbf{Mesafe Seçimi:} Sayısal verilerde \textbf{Euclidean (L2)}, aykırı değerlere hassas durumlarda \textbf{Manhattan (L1)}, metin veya açısal verilerde \textbf{Cosine}, kümeler arası benzerlikte \textbf{Jaccard} kullanılır.
    \item \textbf{Parametre Sayısı (LDA):} LDA, $C$ adet sınıfın olduğu bir problemde en fazla $C-1$ adet yeni boyut (bileşen) üretebilir.
\end{itemize}

\subsection{Karar Ağaçları ve Topluluklar}
\begin{itemize}
    \item \textbf{Safsızlık Ölçütleri:} \textbf{Entropi} (bilgi kazancı odaklı) ve \textbf{Gini Safsızlığı} (daha hızlı hesaplanır, Scikit-Learn varsayılanıdır) ağacı bölmek için kullanılır.
    \item \textbf{Bagging vs. Boosting:} Bagging (Random Forest) varyansı düşürür; Boosting (Gradient Boosting) hem bias hem varyansı düşürerek daha karmaşık hataları çözer.
\end{itemize}

\newpage

% ==============================================================================
% BÖLÜM 2: ÖĞRENME KARTLARI
% ==============================================================================
\section{Öğrenme Kartları (121 Adet)}

\begin{learningcard}
\SoruHeader
$R^2$ (R-Kare) skoru neyi ölçer? \vspace{10pt}
\tcblower
\CevapHeader
Bağımsız değişkenlerin, bağımlı değişkendeki varyansın ne kadarını açıkladığını.
\end{learningcard}

\begin{learningcard}
\SoruHeader
Adam optimizer'ın temel karakteristiği nedir?
\par\vspace{\baselineskip} \tcblower
\CevapHeader
Adaptif öğrenme hızı kullanarak her parametre için güncellemeyi ayarlar.
\end{learningcard}

\begin{learningcard}
\SoruHeader
Açıklanan Varyans Oranı (EVR) formülü nedir?
\par\vspace{\baselineskip} \tcblower
\CevapHeader
$EVR(k) = \frac{\sum_{i=1}^k \lambda_i}{\sum_{i=1}^d \lambda_i}$
\end{learningcard}

\begin{learningcard}
\SoruHeader
Ağırlıkların hata payına göre güncellenmesini sağlayan matematiksel kural nedir?
\par\vspace{\baselineskip} \tcblower
\CevapHeader
Zincir Kuralı (Chain Rule) tabanlı Geri Yayılım.
\end{learningcard}

\begin{learningcard}
\SoruHeader
Batch Normalization'ın temel amacı nedir?
\par\vspace{\baselineskip} \tcblower
\CevapHeader
Eğitimi hızlandırmak ve gradyan sorunlarını azaltmak için girdi değerlerini normalize etmek.
\end{learningcard}

\begin{learningcard}
\SoruHeader
Bias (Önyargı) ve Varyans arasındaki ödünleşim (trade-off) neyi ifade eder?
\par\vspace{\baselineskip} \tcblower
\CevapHeader
Model karmaşıklığı arttıkça biasın azalması ancak varyansın artması durumunu.
\end{learningcard}

\begin{learningcard}
\SoruHeader
Bias-Variance ödünleşimi açısından 'Underfitting' (Eksik Öğrenme) ne zaman gerçekleşir?
\par\vspace{\baselineskip} \tcblower
\CevapHeader
Model çok basit olduğunda ve verideki temel yapıyı öğrenemediğinde (Yüksek Bias).
\end{learningcard}

\begin{learningcard}
\SoruHeader
Bilgi Kazancı (Information Gain) hesaplanırken hangi kavram temel alınır?
\par\vspace{\baselineskip} \tcblower
\CevapHeader
Entropi (H).
\end{learningcard}

\begin{learningcard}
\SoruHeader
Boyut düşürme neden 'Gürültü Azaltma' (Noise Reduction) sağlar?
\par\vspace{\baselineskip} \tcblower
\CevapHeader
Az varyans taşıyan bileşenler (genellikle gürültü) atıldığı için sinyal/gürültü oranı artar.
\end{learningcard}

\begin{learningcard}
\SoruHeader
Boyut düşürme yöntemlerinden PCA (Temel Bileşen Analizi) hangi tür öğrenme kategorisine girer?
\par\vspace{\baselineskip} \tcblower
\CevapHeader
Denetimsiz (Unsupervised) Öğrenme.
\end{learningcard}

\begin{learningcard}
\SoruHeader
Boyut düşürmede 'Kernel PCA' hangi tür veriler için uygundur?
\par\vspace{\baselineskip} \tcblower
\CevapHeader
Doğrusal olarak ayrılamayan, karmaşık yapılara sahip veriler için.
\end{learningcard}

\begin{learningcard}
\SoruHeader
Boyut düşürmede 'Özellik Seçme' (Feature Selection) ile 'Özellik Çıkarma' arasındaki temel fark nedir?
\par\vspace{\baselineskip} \tcblower
\CevapHeader
Özellik seçme orijinal değişkenlerin bir alt kümesini korurken, özellik çıkarma veriyi yeni yapay bileşenlere dönüştürür.
\end{learningcard}

\begin{learningcard}
\SoruHeader
Boyutların Laneti (Curse of Dimensionality) nedeniyle yüksek boyutlarda veri noktaları arasındaki mesafe nasıl değişir?
\par\vspace{\baselineskip} \tcblower
\CevapHeader
Tüm noktalar birbirine neredeyse eşit uzaklıkta görünür ve 'yakınlık' kavramı anlamını yitirir.
\end{learningcard}

\begin{learningcard}
\SoruHeader
Boyutların Laneti (Curse of Dimensionality) veri noktaları arasındaki mesafeyi nasıl etkiler?
\par\vspace{\baselineskip} \tcblower
\CevapHeader
Boyut arttıkça noktalar arası mesafe birbirine çok benzer hale gelir ve en yakın komşu kavramı anlamını yitirir.
\end{learningcard}

\begin{learningcard}
\SoruHeader
DBSCAN algoritması geleneksel merkez tabanlı algoritmalardan hangi özelliği ile ayrılır?
\par\vspace{\baselineskip} \tcblower
\CevapHeader
Yoğunluk farklarına dayanarak karmaşık şekilli kümeleri bulabilmesi.
\end{learningcard}

\begin{learningcard}
\SoruHeader
DBSCAN gibi yoğunluk tabanlı kümeleme algoritmalarının klasik yöntemlere göre en büyük avantajı nedir?
\par\vspace{\baselineskip} \tcblower
\CevapHeader
Düzensiz ve karmaşık şekilli kümeleri tespit edebilmeleri ve gürültüyü (aykırı değerleri) ayırt edebilmeleridir.
\end{learningcard}

\begin{learningcard}
\SoruHeader
Destek Vektör Makineleri (SVM) algoritmasının temel optimizasyon hedefi nedir?
\par\vspace{\baselineskip} \tcblower
\CevapHeader
Sınıflar arasındaki boşluğu (margin) maksimize eden karar hiperdüzlemini bulmak.
\end{learningcard}

\begin{learningcard}
\SoruHeader
Destek Vektör Makineleri'nde (SVM) 'Margin' neyi ifade eder?
\par\vspace{\baselineskip} \tcblower
\CevapHeader
Karar hiperdüzlemi ile en yakın veri noktaları arasındaki mesafe.
\end{learningcard}

\begin{learningcard}
\SoruHeader
Doğrusal Diskriminant Analizi'nin (LDA) temel amacı nedir?
\par\vspace{\baselineskip} \tcblower
\CevapHeader
Sınıflar arası ayrımı (separability) maksimize etmek.
\end{learningcard}

\begin{learningcard}
\SoruHeader
Doğrusal olmayan yapıdaki veriler için PCA'nın hangi türevi 'Kernel Trick' kullanarak çalışır?
\par\vspace{\baselineskip} \tcblower
\CevapHeader
Kernel PCA (Çekirdek Tabanlı PCA).
\end{learningcard}

\begin{learningcard}
\SoruHeader
Doğrusal regresyon veya lojistik regresyonun çözemediği klasik mantıksal problem hangisidir?
\par\vspace{\baselineskip} \tcblower
\CevapHeader
XOR Paradoksu.
\end{learningcard}

\begin{learningcard}
\SoruHeader
Düzeltilmiş R-Kare (Adjusted $R^2$) neden standart $R^2$'den daha güvenilirdir?
\par\vspace{\baselineskip} \tcblower
\CevapHeader
Modele eklenen alakasız değişkenler için penalizasyon uygulayarak 'yalancı başarıyı' engellediği için.
\end{learningcard}

\begin{learningcard}
\SoruHeader
EM (Expectation-Maximization) algoritması kümeleri nasıl tanımlar?
\par\vspace{\baselineskip} \tcblower
\CevapHeader
Veriyi farklı olasılık dağılımlarının (genellikle Gauss) bir bileşimi olarak varsayar ve 'yumuşak' (soft) kümeleme yapar.
\end{learningcard}

\begin{learningcard}
\SoruHeader
Evrensel Yaklaşım Teoremi (Universal Approximation Theorem) neyi savunur?
\par\vspace{\baselineskip} \tcblower
\CevapHeader
Yeterince geniş tek bir gizli katmana sahip MLP'nin her türlü sürekli fonksiyonu temsil edebileceğini.
\end{learningcard}

\begin{learningcard}
\SoruHeader
Geri Yayılım (Backpropagation) algoritmasının temel amacı nedir?
\par\vspace{\baselineskip} \tcblower
\CevapHeader
Çıkış katmanındaki hatayı zincir kuralı kullanarak geriye doğru dağıtmak ve ağırlıkları güncellemek.
\end{learningcard}

\begin{learningcard}
\SoruHeader
Gini Safsızlığı formülü nedir?
\par\vspace{\baselineskip} \tcblower
\CevapHeader
$Gini(t) = 1 - \sum_k p_k^2$
\end{learningcard}

\begin{learningcard}
\SoruHeader
Gradient Boosting yönteminin Random Forest'tan temel farkı nedir?
\par\vspace{\baselineskip} \tcblower
\CevapHeader
Ağaçların paralel değil, bir önceki ağacın hatalarını (artıkları) öğrenecek şekilde ardışıl olarak eğitilmesi.
\end{learningcard}

\begin{learningcard}
\SoruHeader
Gradient Descent algoritmasındaki 'Öğrenme Oranı' ($\alpha$) neyi belirler?
\par\vspace{\baselineskip} \tcblower
\CevapHeader
Minimuma ulaşmak için her adımda atılan güncellemelerin büyüklüğünü.
\end{learningcard}

\begin{learningcard}
\SoruHeader
Gradyan İnişi (Gradient Descent) algoritmasında öğrenme hızı ($\alpha$) neyi belirler?
\par\vspace{\baselineskip} \tcblower
\CevapHeader
Parametre güncellemeleri sırasında atılan adımların büyüklüğünü.
\end{learningcard}

\begin{learningcard}
\SoruHeader
Hiyerarşik kümelemede 'Complete Linkage' mesafeyi nasıl ölçer?
\par\vspace{\baselineskip} \tcblower
\CevapHeader
İki kümenin elemanları arasındaki en büyük mesafeyi temel alarak.
\end{learningcard}

\begin{learningcard}
\SoruHeader
Hiyerarşik kümelemede 'Dendrogram' neyi temsil eder?
\par\vspace{\baselineskip} \tcblower
\CevapHeader
Veri noktaları arasındaki birleşme hiyerarşisini ve kümeler arasındaki mesafe ilişkilerini.
\end{learningcard}

\begin{learningcard}
\SoruHeader
Hiyerarşik kümelemede 'Single Linkage' (Tekil Bağlantı) mesafeyi nasıl tanımlar?
\par\vspace{\baselineskip} \tcblower
\CevapHeader
İki küme arasındaki en yakın eleman çiftleri arasındaki mesafe.
\end{learningcard}

\begin{learningcard}
\SoruHeader
Hiyerarşik kümelemede 'Single Linkage' (Tekli Bağlantı) kuralı nedir?
\par\vspace{\baselineskip} \tcblower
\CevapHeader
İki küme arasındaki mesafeyi, kümelerin birbirine en yakın eleman çifti arasındaki mesafe olarak tanımlar.
\end{learningcard}

\begin{learningcard}
\SoruHeader
İki nokta arasındaki en kısa doğrusal mesafeyi ölçen metriğe ne denir?
\par\vspace{\baselineskip} \tcblower
\CevapHeader
Öklid (Euclidean) Mesafesi.
\end{learningcard}

\begin{learningcard}
\SoruHeader
İkili sınıflandırma (Binary Classification) için kullanılan temel kayıp fonksiyonu (loss function) nedir?
\par\vspace{\baselineskip} \tcblower
\CevapHeader
Binary Cross-Entropy Loss.
\end{learningcard}

\begin{learningcard}
\SoruHeader
k-En Yakın Komşu (k-NN) algoritması neden 'Tembel Öğrenme' (Lazy Learning) olarak adlandırılır?
\par\vspace{\baselineskip} \tcblower
\CevapHeader
Eğitim aşamasında bir model kurmayıp, tüm hesaplamayı tahmin anında yaptığı için.
\end{learningcard}

\begin{learningcard}
\SoruHeader
K-Katlı Çapraz Doğrulama (K-Fold Cross Validation) yönteminin amacı nedir?
\par\vspace{\baselineskip} \tcblower
\CevapHeader
Veriyi farklı parçalara bölerek modelin genelleme yeteneğini daha güvenilir bir şekilde ölçmek.
\end{learningcard}

\begin{learningcard}
\SoruHeader
k-NN'de $K$ sayısının çok küçük seçilmesi hangi soruna yol açar?
\par\vspace{\baselineskip} \tcblower
\CevapHeader
Yüksek varyans ve aşırı öğrenme (overfitting).
\end{learningcard}

\begin{learningcard}
\SoruHeader
K-Ortalamalar (K-Means) algoritmasında 'Inertia' metriği neyi ifade eder?
\par\vspace{\baselineskip} \tcblower
\CevapHeader
Veri noktalarının kendi küme merkezlerine olan uzaklıklarının kareleri toplamını.
\end{learningcard}

\begin{learningcard}
\SoruHeader
K-Ortalamalar (K-Means) algoritmasının optimizasyon hedefi nedir?
\par\vspace{\baselineskip} \tcblower
\CevapHeader
Küme içi varyasyonu (Inertia) minimize etmek.
\end{learningcard}

\begin{learningcard}
\SoruHeader
Karar ağaçlarında 'Dallanma Kriteri' olarak ne seçilir?
\par\vspace{\baselineskip} \tcblower
\CevapHeader
En yüksek bilgi kazancını (Information Gain) sağlayan özellik ve eşik değeri.
\end{learningcard}

\begin{learningcard}
\SoruHeader
Karar Ağaçlarında 'Gini Safsızlığı' (Gini Impurity) neyi ölçer?
\par\vspace{\baselineskip} \tcblower
\CevapHeader
Bir düğümdeki verilerin sınıfsal homojenliğini (karışıklığını).
\end{learningcard}

\begin{learningcard}
\SoruHeader
Karar ağaçlarında dallanma kriteri olarak kullanılan saflık ölçüsü nedir?
\par\vspace{\baselineskip} \tcblower
\CevapHeader
Gini Safsızlığı veya Entropi.
\end{learningcard}

\begin{learningcard}
\SoruHeader
Karmaşıklık Matrisinde (Confusion Matrix) 'F1 Skoru' nasıl hesaplanır?
\par\vspace{\baselineskip} \tcblower
\CevapHeader
Hassasiyet (Precision) ve Duyarlılığın (Recall) harmonik ortalaması alınarak.
\end{learningcard}

\begin{learningcard}
\SoruHeader
KD-Ağaçları (KD-Trees) uzayı bölmek için hangi geometrik yapıyı kullanır?
\par\vspace{\baselineskip} \tcblower
\CevapHeader
Eksenlere dik hiperdüzlemler.
\end{learningcard}

\begin{learningcard}
\SoruHeader
KD-Ağaçları hangi boyut aralığında verimli arama yapabilmektedir?
\par\vspace{\baselineskip} \tcblower
\CevapHeader
Genellikle 2 ile 8 boyut arasında.
\end{learningcard}

\begin{learningcard}
\SoruHeader
KNN algoritmasında 'K' değerinin çok küçük seçilmesi hangi soruna yol açar?
\par\vspace{\baselineskip} \tcblower
\CevapHeader
Aşırı öğrenme (Overfitting/Yüksek Varyans).
\end{learningcard}

\begin{learningcard}
\SoruHeader
Kök Ortalama Kare Hata (RMSE) metriğinin MAE'ye göre yorumlama avantajı nedir?
\par\vspace{\baselineskip} \tcblower
\CevapHeader
Hata birimini orijinal hedef değişken birimine (Örn: Dolar, Derece) geri döndürmesi.
\end{learningcard}

\begin{learningcard}
\SoruHeader
Kümeleme (Clustering) analizinin temel amacı nedir?
\par\vspace{\baselineskip} \tcblower
\CevapHeader
Grup içi benzerliği en yüksek, gruplar arası farklılığı en fazla olacak şekilde veriyi alt gruplara ayırmak.
\end{learningcard}

\begin{learningcard}
\SoruHeader
Kümeleme Analizi (Clustering) hangi tür öğrenme yöntemidir?
\par\vspace{\baselineskip} \tcblower
\CevapHeader
Gözetimsiz (Unsupervised) Öğrenme.
\end{learningcard}

\begin{learningcard}
\SoruHeader
Kümeleme analizinde 'Öklid Mesafesi' nasıl tanımlanır?
\par\vspace{\baselineskip} \tcblower
\CevapHeader
İki nokta arasındaki en kısa, doğrusal mesafe.
\end{learningcard}

\begin{learningcard}
\SoruHeader
Kümeleme sonuçlarını ağaç benzeri bir yapıda gösteren diyagrama ne ad verilir?
\par\vspace{\baselineskip} \tcblower
\CevapHeader
Dendrogram.
\end{learningcard}

\begin{learningcard}
\SoruHeader
Kümelemede 'Uzman Doğrulaması' (Sanity Check) neden gereklidir?
\par\vspace{\baselineskip} \tcblower
\CevapHeader
İstatistiksel olarak bulunan kümelerin gerçek dünya iş mantığına (Örn: Pazar segmentasyonu) uygunluğunu denetlemek için.
\end{learningcard}

\begin{learningcard}
\SoruHeader
L1 Düzenlileştirmesi (Lasso) modelde nasıl bir etki yaratır?
\par\vspace{\baselineskip} \tcblower
\CevapHeader
Bazı katsayıları tam olarak sıfıra eşitleyerek özellik seçimi (feature selection) yapar.
\end{learningcard}

\begin{learningcard}
\SoruHeader
L2 Düzenlileştirmesi (Ridge) maliyet fonksiyonuna ne ekler?
\par\vspace{\baselineskip} \tcblower
\CevapHeader
Ağırlıkların karelerinin toplamını ($\lambda \cdot \sum w_j^2$).
\end{learningcard}

\begin{learningcard}
\SoruHeader
LDA (Doğrusal Diskriminant Analizi) ile PCA arasındaki en temel kavramsal fark nedir?
\par\vspace{\baselineskip} \tcblower
\CevapHeader
PCA denetimsiz (etiketsiz) bir yöntemken, LDA sınıfları birbirinden ayırmak için etiketli veri kullanan denetimli bir yöntemdir.
\end{learningcard}

\begin{learningcard}
\SoruHeader
LDA'da 'Sınıf İçi Saçılım' ($S_W$) matrisinin minimize edilmesi ne anlama gelir?
\par\vspace{\baselineskip} \tcblower
\CevapHeader
Aynı sınıfa ait örneklerin birbirine olabildiğince yakınlaşması.
\end{learningcard}

\begin{learningcard}
\SoruHeader
LDA'da maksimum bileşen sayısı ne ile sınırlıdır?
\par\vspace{\baselineskip} \tcblower
\CevapHeader
Sınıf sayısı $K$ olmak üzere, en fazla $K-1$ bileşen üretilebilir.
\end{learningcard}

\begin{learningcard}
\SoruHeader
LDA'da üretilebilecek maksimum bileşen sayısı ne ile sınırlıdır?
\par\vspace{\baselineskip} \tcblower
\CevapHeader
Sınıf sayısı eksi bir ($K-1$).
\end{learningcard}

\begin{learningcard}
\SoruHeader
Lojistik Regresyon modelinde çıktı olasılığını $[0, 1]$ aralığına sıkıştıran fonksiyon hangisidir?
\par\vspace{\baselineskip} \tcblower
\CevapHeader
Sigmoid (Lojistik) fonksiyonu: $\sigma(z) = \frac{1}{1 + e^{-z}}$.
\end{learningcard}

\begin{learningcard}
\SoruHeader
Lojistik regresyonda katsayıların aşırı büyümesini engellemek için kullanılan L2 düzenlileştirmesinin diğer adı nedir?
\par\vspace{\baselineskip} \tcblower
\CevapHeader
Ridge Düzenlileştirme.
\end{learningcard}

\begin{learningcard}
\SoruHeader
Manhattan (City-block) mesafesi hangi kurala göre hesaplanır?
\par\vspace{\baselineskip} \tcblower
\CevapHeader
Koordinat eksenleri boyunca hareket edilerek hesaplanan mutlak farkların toplamı.
\end{learningcard}

\begin{learningcard}
\SoruHeader
Manhattan mesafesi nasıl hesaplanır?
\par\vspace{\baselineskip} \tcblower
\CevapHeader
Koordinat eksenleri boyunca mutlak farkların toplamı ile.
\end{learningcard}

\begin{learningcard}
\SoruHeader
MAPE metriği hangi durumlarda hatalı veya tanımsız sonuçlar üretir?
\par\vspace{\baselineskip} \tcblower
\CevapHeader
Gerçek değerler sıfır veya sıfıra çok yakın olduğunda.
\end{learningcard}

\begin{learningcard}
\SoruHeader
Mevcut özelliklerden bir alt küme seçerek orijinal değişkenlerin korunduğu işleme ne denir?
\par\vspace{\baselineskip} \tcblower
\CevapHeader
Özellik Seçme (Feature Selection).
\end{learningcard}

\begin{learningcard}
\SoruHeader
Modelin eğitim verisine aşırı uyum sağlaması sonucu test verisinde düşük performans göstermesi durumuna ne denir?
\par\vspace{\baselineskip} \tcblower
\CevapHeader
Overfitting (Aşırı Öğrenme).
\end{learningcard}

\begin{learningcard}
\SoruHeader
Naive Bayes algoritmasının 'Naive' (Saf) olarak adlandırılmasının sebebi nedir?
\par\vspace{\baselineskip} \tcblower
\CevapHeader
Özelliklerin sınıf verildiğinde birbirinden tamamen bağımsız olduğunu varsaymasıdır.
\end{learningcard}

\begin{learningcard}
\SoruHeader
Naive Bayes türlerinden hangisi sürekli (continuous) değişkenler için uygundur?
\par\vspace{\baselineskip} \tcblower
\CevapHeader
Gaussian Naive Bayes.
\end{learningcard}

\begin{learningcard}
\SoruHeader
Neden 'Düzeltilmiş R-Kare' (Adjusted $R^2$) metriği standart R-Kare'den daha güvenilirdir?
\par\vspace{\baselineskip} \tcblower
\CevapHeader
Modele eklenen anlamsız değişkenler için cezalandırma yaparak sahte başarıyı engellediği için.
\end{learningcard}

\begin{learningcard}
\SoruHeader
Neden PCA uygulanmadan önce verinin standartlaştırılması (z-score scaling) önerilir?
\par\vspace{\baselineskip} \tcblower
\CevapHeader
Ölçeği büyük olan değişkenlerin varyans hesaplamasını domine etmesini engellemek için.
\end{learningcard}

\begin{learningcard}
\SoruHeader
Ortalama Kare Hata (MSE) metriği hataları nasıl cezalandırır?
\par\vspace{\baselineskip} \tcblower
\CevapHeader
Hataların karesini aldığı için büyük sapmaları çok daha sert (üstel olarak) cezalandırır.
\end{learningcard}

\begin{learningcard}
\SoruHeader
Ortalama Kare Hata (MSE) metriğinin MAE'den temel farkı nedir?
\par\vspace{\baselineskip} \tcblower
\CevapHeader
Hataların karesini aldığı için büyük hataları daha sert cezalandırması.
\end{learningcard}

\begin{learningcard}
\SoruHeader
Ortalama Mutlak Hata (MAE) metriğinin en güçlü yönü nedir?
\par\vspace{\baselineskip} \tcblower
\CevapHeader
Aykırı değerlere (outliers) karşı dirençli ve stabil bir ölçüm sunması.
\end{learningcard}

\begin{learningcard}
\SoruHeader
Ortalama Mutlak Yüzde Hata (MAPE) metriğinin temel kısıtlaması nedir?
\par\vspace{\baselineskip} \tcblower
\CevapHeader
Gerçek değerler sıfıra çok yakın olduğunda sonucun matematiksel olarak belirsizleşmesi veya patlaması.
\end{learningcard}

\begin{learningcard}
\SoruHeader
Otoenkodör türlerinden hangisi gürültülü veriden temiz veriyi öğrenmek için tasarlanmıştır?
\par\vspace{\baselineskip} \tcblower
\CevapHeader
Denoising Autoencoder.
\end{learningcard}

\begin{learningcard}
\SoruHeader
Otoenkodörlerde (Autoencoders) 'Bottleneck' katmanının görevi nedir?
\par\vspace{\baselineskip} \tcblower
\CevapHeader
Veriyi en düşük boyutlu (sıkıştırılmış) temsiline ($z$) indirgemek.
\end{learningcard}

\begin{learningcard}
\SoruHeader
Otoenkodörlerde (Autoencoders) verinin en düşük boyutlu temsilinin bulunduğu katmana ne denir?
\par\vspace{\baselineskip} \tcblower
\CevapHeader
Bottleneck (Darboğaz).
\end{learningcard}

\begin{learningcard}
\SoruHeader
PCA (Temel Bileşen Analizi) algoritmasının ilk adımı nedir?
\par\vspace{\baselineskip} \tcblower
\CevapHeader
Veriyi merkezlemek (sütun ortalamasını her gözlemden çıkarmak).
\end{learningcard}

\begin{learningcard}
\SoruHeader
PCA algoritmasında özdeğerlerin ($\lambda_i$) büyüklüğü neyi temsil eder?
\par\vspace{\baselineskip} \tcblower
\CevapHeader
İlgili bileşenin taşıdığı varyans miktarını.
\end{learningcard}

\begin{learningcard}
\SoruHeader
PCA formülasyonunda özvektörler arasındaki ilişki nasıldır?
\par\vspace{\baselineskip} \tcblower
\CevapHeader
Özvektörler birbirine ortogonaldir (diktir) ve aralarında korelasyon yoktur.
\end{learningcard}

\begin{learningcard}
\SoruHeader
PCA'da açıklanan varyans oranı (EVR) nasıl hesaplanır?
\par\vspace{\baselineskip} \tcblower
\CevapHeader
$EVR(k) = \frac{\sum_{i=1}^k \lambda_i}{\sum_{i=1}^d \lambda_i}$
\end{learningcard}

\begin{learningcard}
\SoruHeader
PCA'da en uygun bileşen sayısını ($k$) belirlemek için kullanılan görsel grafiğe ne ad verilir?
\par\vspace{\baselineskip} \tcblower
\CevapHeader
Scree Plot (Yamaç Grafiği).
\end{learningcard}

\begin{learningcard}
\SoruHeader
PCA'da her bir ana bileşenin (PC) bir sonrakine göre geometrik konumu nedir?
\par\vspace{\baselineskip} \tcblower
\CevapHeader
Birbirlerine diktirler (Ortogonal).
\end{learningcard}

\begin{learningcard}
\SoruHeader
PCA'da kovaryans matrisi $C$ hangi matematiksel formülle hesaplanır?
\par\vspace{\baselineskip} \tcblower
\CevapHeader
$C = \frac{1}{N} \tilde{X}^T \tilde{X}$ (Burada $\tilde{X}$ merkezlenmiş veridir).
\end{learningcard}

\begin{learningcard}
\SoruHeader
PCA'da kullanılan 'Scree Plot' grafiği neyi belirlemek için kullanılır?
\par\vspace{\baselineskip} \tcblower
\CevapHeader
Özdeğerlerin sıralı dizilimindeki 'dirsek' noktasını bularak optimal bileşen sayısını ($k$) seçmek.
\end{learningcard}

\begin{learningcard}
\SoruHeader
PCA'da verinin merkezlenmesi için hangi işlem yapılır?
\par\vspace{\baselineskip} \tcblower
\CevapHeader
$\tilde{X} = X - \mu$ (Her özellikten sütun ortalamasının çıkarılması).
\end{learningcard}

\begin{learningcard}
\SoruHeader
PCA'nın geometrik yorumunda birinci temel bileşen (PC1) hangi yöndedir?
\par\vspace{\baselineskip} \tcblower
\CevapHeader
Verideki maksimum varyansın olduğu doğrultudadır.
\end{learningcard}

\begin{learningcard}
\SoruHeader
R-Kare ($R^2$) skoru neyi ifade eder?
\par\vspace{\baselineskip} \tcblower
\CevapHeader
Hedef değişkendeki varyansın yüzde kaçının modeldeki özellikler tarafından açıklandığını.
\end{learningcard}

\begin{learningcard}
\SoruHeader
Random Forest'ta her düğümde rastgele $m < d$ özellik seçilmesinin temel amacı nedir?
\par\vspace{\baselineskip} \tcblower
\CevapHeader
Ağaçlar arasındaki korelasyonu düşürerek topluluk varyansını azaltmak.
\end{learningcard}

\begin{learningcard}
\SoruHeader
Rastgele Orman (Random Forest) algoritmasındaki 'Bagging' mantığı nedir?
\par\vspace{\baselineskip} \tcblower
\CevapHeader
Eğitim verisinden tekrarlı rastgele örneklemler (bootstrap) alarak birden fazla ağaç eğitmek ve sonuçları oylamak.
\end{learningcard}

\begin{learningcard}
\SoruHeader
Rastgele Orman (Random Forest) algoritmasının temel çalışma prensibi nedir?
\par\vspace{\baselineskip} \tcblower
\CevapHeader
Birden fazla karar ağacını bootstrap örnekleme ve rastgele özellik seçimiyle birleştirmek.
\end{learningcard}

\begin{learningcard}
\SoruHeader
Regresyon modelinde $R^2 = 0.82$ sonucu ne anlama gelir?
\par\vspace{\baselineskip} \tcblower
\CevapHeader
Model, verideki değişkenliğin $\%82$'sini açıklayabilmektedir.
\end{learningcard}

\begin{learningcard}
\SoruHeader
Regresyon modellerinde 'Artık Değer' (Residual) nedir?
\par\vspace{\baselineskip} \tcblower
\CevapHeader
Gerçek değer ($y_i$) ile tahmin edilen değer ($\hat{y}_i$) arasındaki fark.
\end{learningcard}

\begin{learningcard}
\SoruHeader
RMSE metriği MSE'ye göre neden daha yorumlanabilirdir?
\par\vspace{\baselineskip} \tcblower
\CevapHeader
Karekök işlemi sayesinde hata birimini orijinal verinin birimine döndürdüğü için.
\end{learningcard}

\begin{learningcard}
\SoruHeader
Sigmoid fonksiyonunun türevinin çok küçük olması hangi eğitim problemine yol açar?
\par\vspace{\baselineskip} \tcblower
\CevapHeader
Gradyan Kaybolması (Vanishing Gradient).
\end{learningcard}

\begin{learningcard}
\SoruHeader
Sinir ağlarında 'Batch Normalization' neden kullanılır?
\par\vspace{\baselineskip} \tcblower
\CevapHeader
Eğitimi hızlandırmak ve gradyan sorunlarını (vanishing gradient) azaltmak için katman girişlerini normalize eder.
\end{learningcard}

\begin{learningcard}
\SoruHeader
Sinir ağlarında 'Dropout' tekniği neden kullanılır?
\par\vspace{\baselineskip} \tcblower
\CevapHeader
Nöronları rastgele kapatarak aşırı öğrenmeyi (overfitting) engellemek için.
\end{learningcard}

\begin{learningcard}
\SoruHeader
Sinir ağlarında 'ReLU' aktivasyon fonksiyonunun tanımı nedir?
\par\vspace{\baselineskip} \tcblower
\CevapHeader
$g(z) = \max(0, z)$.
\end{learningcard}

\begin{learningcard}
\SoruHeader
Sinir ağlarında 'Softmax' fonksiyonu genellikle hangi katmanda kullanılır?
\par\vspace{\baselineskip} \tcblower
\CevapHeader
Çok sınıflı sınıflandırma problemlerinde çıkış katmanında.
\end{learningcard}

\begin{learningcard}
\SoruHeader
Sinir ağlarında katmanlar boyunca verinin işlenerek tahmin üretilmesi sürecine ne ad verilir?
\par\vspace{\baselineskip} \tcblower
\CevapHeader
İleri Besleme (Forward Propagation).
\end{learningcard}

\begin{learningcard}
\SoruHeader
Sinir ağlarında negatif değerleri sıfıra eşitleyen en yaygın aktivasyon fonksiyonu hangisidir?
\par\vspace{\baselineskip} \tcblower
\CevapHeader
ReLU ($g(z) = \max(0, z)$).
\end{learningcard}

\begin{learningcard}
\SoruHeader
SOM (Kendi Kendini Organize Eden Haritalar) sinir ağı mimarisinin temel işlevi nedir?
\par\vspace{\baselineskip} \tcblower
\CevapHeader
Yüksek boyutlu veriyi topolojiyi koruyarak düşük boyutlu (genellikle 2D) bir haritaya indirgemek.
\end{learningcard}

\begin{learningcard}
\SoruHeader
SVD (Tekil Değer Ayrışımı) ile PCA arasındaki ilişki nedir?
\par\vspace{\baselineskip} \tcblower
\CevapHeader
SVD, PCA hesaplamasını sayısal olarak daha kararlı bir şekilde gerçekleştirmek için kullanılır.
\end{learningcard}

\begin{learningcard}
\SoruHeader
SVM algoritmasında 'Soft Margin' kullanımı neyi amaçlar?
\par\vspace{\baselineskip} \tcblower
\CevapHeader
Verideki bazı gürültülü noktaların yanlış sınıfta kalmasına izin vererek daha genel bir karar sınırı oluşturmayı.
\end{learningcard}

\begin{learningcard}
\SoruHeader
SVM'de 'Kernel Trick' ne işe yarar?
\par\vspace{\baselineskip} \tcblower
\CevapHeader
Veriyi daha yüksek boyutlu bir özellik uzayına taşıyarak doğrusal olmayan sınırların çizilmesini sağlar.
\end{learningcard}

\begin{learningcard}
\SoruHeader
SVM'de veriyi doğrusal olmayan bir üst uzaya taşıyarak ayrıştırmayı sağlayan yönteme ne denir?
\par\vspace{\baselineskip} \tcblower
\CevapHeader
Kernel Trick (Çekirdek Yöntemi).
\end{learningcard}

\begin{learningcard}
\SoruHeader
Sınıflandırma (Classification) problemi makine öğrenmesinin hangi kategorisine girer?
\par\vspace{\baselineskip} \tcblower
\CevapHeader
Denetimli Öğrenme (Supervised Learning).
\end{learningcard}

\begin{learningcard}
\SoruHeader
Sınıflandırma performansını ölçen Karmaşıklık Matrisi'nde (Confusion Matrix) FN neyi temsil eder?
\par\vspace{\baselineskip} \tcblower
\CevapHeader
Yanlış Negatif (Gerçekte pozitif olanın negatif tahmin edilmesi).
\end{learningcard}

\begin{learningcard}
\SoruHeader
Sınıflandırma probleminde modelin çıktı olarak ürettiği olasılığı 0 veya 1'e dönüştüren fonksiyona ne denir?
\par\vspace{\baselineskip} \tcblower
\CevapHeader
Sigmoid Fonksiyonu: $\sigma(z) = \frac{1}{1 + e^{-z}}$.
\end{learningcard}

\begin{learningcard}
\SoruHeader
t-SNE algoritması yüksek boyuttaki benzerlikleri düşük boyuta taşırken hangi olasılık dağılımını kullanır?
\par\vspace{\baselineskip} \tcblower
\CevapHeader
t-dağılımı (Student's t-distribution).
\end{learningcard}

\begin{learningcard}
\SoruHeader
UMAP algoritmasının t-SNE'ye göre en büyük pratik avantajı nedir?
\par\vspace{\baselineskip} \tcblower
\CevapHeader
Daha hızlı çalışması ve yeni veri noktalarının projeksiyonuna izin vermesi.
\end{learningcard}

\begin{learningcard}
\SoruHeader
UMAP yönteminin t-SNE'ye göre avantajı nedir?
\par\vspace{\baselineskip} \tcblower
\CevapHeader
Daha hızlı çalışması ve yeni veri noktalarını mevcut projeksiyona dahil edebilmesidir.
\end{learningcard}

\begin{learningcard}
\SoruHeader
Vektörleştirme (Vectorization) sinir ağı eğitiminde neden önemlidir?
\par\vspace{\baselineskip} \tcblower
\CevapHeader
Çok sayıda skaler işlemi tek bir matris çarpımına dönüştürerek donanım (GPU) hızlandırması sağlar.
\end{learningcard}

\begin{learningcard}
\SoruHeader
Veriyi 2D veya 3D'de görselleştirmek için t-SNE'nin kullandığı temel benzerlik ölçütü nedir?
\par\vspace{\baselineskip} \tcblower
\CevapHeader
Koşullu olasılıklar tabanlı benzerlik (t-dağılımı).
\end{learningcard}

\begin{learningcard}
\SoruHeader
XOR problemi neden tek katmanlı (doğrusal) modellerle çözülemez?
\par\vspace{\baselineskip} \tcblower
\CevapHeader
XOR verisi doğrusal olarak ayrıştırılamaz, modellemek için doğrusal olmayan bir gizli katman gerekir.
\end{learningcard}

\begin{learningcard}
\SoruHeader
Yapay sinir ağlarında 'Dropout' yönteminin işlevi nedir?
\par\vspace{\baselineskip} \tcblower
\CevapHeader
Eğitim sırasında nöronları rastgele kapatarak modelin belirli nöronlara aşırı bağımlı olmasını (overfitting) engellemek.
\end{learningcard}

\begin{learningcard}
\SoruHeader
Yapay sinir ağlarında bir nöronun temel matematiksel denklemi nedir?
\par\vspace{\baselineskip} \tcblower
\CevapHeader
$a = \phi(\sum w_j x_j + b)$ (Burada $\phi$ aktivasyon fonksiyonudur).
\end{learningcard}

\begin{learningcard}
\SoruHeader
Yapay sinir ağlarında tek bir nöronun matematiksel denklemi nedir?
\par\vspace{\baselineskip} \tcblower
\CevapHeader
$a = \phi(\sum w_j x_j + b)$
\end{learningcard}

\begin{learningcard}
\SoruHeader
Yüksek boyutlu veriyi 2D veya 3D uzayda görselleştirmek için en sık tercih edilen doğrusal olmayan yöntemler hangileridir?
\par\vspace{\baselineskip} \tcblower
\CevapHeader
t-SNE ve UMAP yöntemleri.
\end{learningcard}

\begin{learningcard}
\SoruHeader
Çok sınıflı sınıflandırma problemlerinde kullanılan hata fonksiyonu hangisidir?
\par\vspace{\baselineskip} \tcblower
\CevapHeader
Çapraz Entropi (Cross-Entropy Loss).
\end{learningcard}

\begin{learningcard}
\SoruHeader
Özelliklerin dönüştürülerek yeni yapay bileşenlerin üretildiği işleme ne ad verilir?
\par\vspace{\baselineskip} \tcblower
\CevapHeader
Boyut Düşürme (Dimensionality Reduction).
\end{learningcard}

\newpage

% ==============================================================================
% BÖLÜM 3: TEST BANKASI
% ==============================================================================
\section{Test Bankası (90 Soru)}

\subsection{Algoritma Testi}

\begin{testcard}{Soru 1}
    Bir veri setinde bağımsız değişkenlerin (X) bağımlı değişkeni (y) tahmin etmek için kullanıldığı, k-NN veya Lojistik Regresyon gibi algoritmaları kapsayan öğrenme türü hangisidir?
    \begin{itemize}
        \item[A)] Denetimsiz Öğrenme
        \item[B)] Denetimli Öğrenme
        \item[C)] Takviyeli Öğrenme
        \item[D)] Yarı Denetimli Öğrenme
    \end{itemize}
    \tcblower
    \textbf{Doğru Cevap: B} \\
    \textbf{Açıklama:} Denetimli öğrenmede her girdinin bir karşılık gelen etiketi bulunur; model bu eşleşmeyi öğrenir.
\end{testcard}

\begin{testcard}{Soru 2}
    Denetimsiz öğrenme algoritmalarının temel amacı aşağıdakilerden hangisidir?
    \begin{itemize}
        \item[A)] Veri setindeki etiketleri tahmin etmek
        \item[B)] Verideki gizli yapıları ve grupları (kümeleri) keşfetmek
        \item[C)] Bir regresyon doğrusu oluşturmak
        \item[D)] Hata payını sıfıra indirmek
    \end{itemize}
    \tcblower
    \textbf{Doğru Cevap: B} \\
    \textbf{Açıklama:} Denetimsiz öğrenmede etiket yoktur, amaç verinin kendi içindeki benzerlikleri bulmaktır.
\end{testcard}

\begin{testcard}{Soru 3}
    "Eğer bir ördek gibi yürüyorsa ve ördek gibi vakvaklıyorsa, o bir ördektir" mantığıyla çalışan, en yakın komşularına bakarak sınıflandırma yapan algoritma hangisidir?
    \begin{itemize}
        \item[A)] k-NN (k-En Yakın Komşu)
        \item[B)] Destek Vektör Makineleri (SVM)
        \item[C)] Karar Ağaçları
        \item[D)] PCA
    \end{itemize}
    \tcblower
    \textbf{Doğru Cevap: A} \\
    \textbf{Açıklama:} k-NN, "benzer şeyler birbirine yakındır" ilkesiyle en yakın komşuların oylamasına dayanır.
\end{testcard}

\begin{testcard}{Soru 4}
    K-Means algoritmasında, verilerin hangi kümeye ait olduğunu belirlemek için genellikle hangi mesafe ölçütü kullanılır?
    \begin{itemize}
        \item[A)] Manhattan Mesafesi
        \item[B)] Öklid Mesafesi
        \item[C)] Jaccard Benzerliği
        \item[D)] Kosinüs Benzerliği
    \end{itemize}
    \tcblower
    \textbf{Doğru Cevap: B} \\
    \textbf{Açıklama:} K-Means, küme içi varyansı minimize etmek için noktaların merkezlere olan Öklid mesafesini kullanır.
\end{testcard}

\begin{testcard}{Soru 5}
    Karar Ağaçlarında (Decision Trees), bir düğümü bölmek için kullanılan ve "safsızlık" (impurity) ölçen temel metriklerden biri hangisidir?
    \begin{itemize}
        \item[A)] Gini Katsayısı / Entropi
        \item[B)] R-Kare
        \item[C)] F1 Skoru
        \item[D)] Ortalama Mutlak Hata (MAE)
    \end{itemize}
    \tcblower
    \textbf{Doğru Cevap: A} \\
    \textbf{Açıklama:} Karar ağaçları, dallanma kararlarını verinin saflığını artıran Gini veya Bilgi Kazancı değerlerine göre verir.
\end{testcard}

\begin{testcard}{Soru 6}
    Destek Vektör Makineleri (SVM) algoritmasında, sınıfları birbirinden ayıran en geniş boşluğu (margin) bulmak için kullanılan veri noktalarına ne ad verilir?
    \begin{itemize}
        \item[A)] Aykırı Değerler (Outliers)
        \item[B)] Destek Vektörleri (Support Vectors)
        \item[C)] Merkez Noktaları (Centroids)
        \item[D)] Karar Sınırları
    \end{itemize}
    \tcblower
    \textbf{Doğru Cevap: B} \\
    \textbf{Açıklama:} Destek vektörleri, karar sınırını "ayakta tutan" ve sınırın yerini belirleyen en yakın örneklerdir.
\end{testcard}

\begin{testcard}{Soru 7}
    Rastgele Orman (Random Forest) algoritması, birden fazla karar ağacını birleştirerek daha güçlü bir model oluşturur. Bu yönteme genel olarak ne ad verilir?
    \begin{itemize}
        \item[A)] Boosting
        \item[B)] Bagging (Bootstrap Aggregating)
        \item[C)] Dimensionality Reduction
        \item[D)] Clustering
    \end{itemize}
    \tcblower
    \textbf{Doğru Cevap: B} \\
    \textbf{Açıklama:} Random Forest, veriden farklı alt kümeler çekerek çok sayıda ağaç eğitir ve sonuçları oylar.
\end{testcard}

\begin{testcard}{Soru 8}
    Bir modelin eğitim verisine çok iyi uyum sağlayıp, test (yeni) verisinde kötü performans göstermesi durumuna ne ad verilir?
    \begin{itemize}
        \item[A)] Eksik Öğrenme (Underfitting)
        \item[B)] Aşırı Öğrenme (Overfitting)
        \item[C)] Doğrusal Regresyon
        \item[D)] Bias (Yanlılık)
    \end{itemize}
    \tcblower
    \textbf{Doğru Cevap: B} \\
    \textbf{Açıklama:} Model verideki gürültüyü ezberlemiş ve yeni verilere karşı genelleme yapamaz hale gelmiştir.
\end{testcard}

\begin{testcard}{Soru 9}
    Yapay Sinir Ağlarında, girdilerin ağırlıklarla çarpılıp toplandığı ve bir eşik değerinden geçirilerek çıktının üretildiği temel yapı birimine ne denir?
    \begin{itemize}
        \item[A)] Katman (Layer)
        \item[B)] Nöron (Perceptron)
        \item[C)] Aktivasyon Fonksiyonu
        \item[D)] Optimizer
    \end{itemize}
    \tcblower
    \textbf{Doğru Cevap: B} \\
    \textbf{Açıklama:} Perceptron, bir sinir ağındaki en küçük ve temel işlem birimidir.
\end{testcard}

\begin{testcard}{Soru 10}
    Boyut indirgeme (Dimensionality Reduction) amacıyla kullanılan ve verideki varyansı en çok koruyan yeni bileşenleri bulan algoritma hangisidir?
    \begin{itemize}
        \item[A)] K-Means
        \item[B)] PCA (Temel Bileşen Analizi)
        \item[C)] Lojistik Regresyon
        \item[D)] Destek Vektör Makineleri (SVM)
    \end{itemize}
    \tcblower
    \textbf{Doğru Cevap: B} \\
    \textbf{Açıklama:} PCA (Principal Component Analysis), verideki en yüksek varyansı koruyacak şekilde boyut indirgeme yapan temel algoritmadır.
\end{testcard}

\begin{testcard}{Soru 11}
    Boyut indirgeme (Dimensionality Reduction) amacıyla kullanılan ve verideki maksimum varyansı koruyarak yeni bileşenler oluşturan algoritma hangisidir?
    \begin{itemize}
        \item[A)] DBSCAN
        \item[B)] PCA (Temel Bileşen Analizi)
        \item[C)] Naive Bayes
        \item[D)] Lojistik Regresyon
    \end{itemize}
    \tcblower
    \textbf{Doğru Cevap: B} \\
    \textbf{Açıklama:} PCA, verideki en yüksek değişkenliği temsil eden ana yönleri (temel bileşenler) bularak bilgiyi sıkıştırır.
\end{testcard}

\begin{testcard}{Soru 12}
    k-NN algoritmasında "k" parametresinin çok küçük seçilmesi (örneğin k=1) aşağıdakilerden hangisine neden olabilir?
    \begin{itemize}
        \item[A)] Modelin çok basitleşmesine
        \item[B)] Aşırı öğrenmeye (Overfitting) ve gürültüye duyarlılığa
        \item[C)] Eksik öğrenmeye (Underfitting)
        \item[D)] Hesaplama maliyetinin azalmasına
    \end{itemize}
    \tcblower
    \textbf{Doğru Cevap: B} \\
    \textbf{Açıklama:} k=1 olduğunda model her noktayı kendisine en yakın tek bir komşuya göre atar, bu da verideki rastgele gürültülerin model sınırlarını bozmasına yol açar.
\end{testcard}

\begin{testcard}{Soru 13}
    Lojistik Regresyon algoritması, doğrusal regresyonun aksine çıktıyı 0 ile 1 arasında bir olasılık değerine dönüştürmek için hangi fonksiyonu kullanır?
    \begin{itemize}
        \item[A)] ReLU Fonksiyonu
        \item[B)] Sigmoid (Lojistik) Fonksiyonu
        \item[C)] Tanh Fonksiyonu
        \item[D)] Step Fonksiyonu
    \end{itemize}
    \tcblower
    \textbf{Doğru Cevap: B} \\
    \textbf{Açıklama:} Sigmoid fonksiyonu, her türlü sayısal değeri [0, 1] aralığına haritalayarak modelin "sınıf olasılığı" üretmesini sağlar.
\end{testcard}

\begin{testcard}{Soru 14}
    Naive Bayes algoritmasının ismindeki "Naive" (Saf/Safdil) ifadesi hangi varsayımdan kaynaklanır?
    \begin{itemize}
        \item[A)] Verilerin doğrusal olduğu
        \item[B)] Özelliklerin (features) birbirinden bağımsız olduğu varsayım
        \item[C)] Veride hiç gürültü olmadığı
        \item[D)] Tüm verilerin normal dağıldığı
    \end{itemize}
    \tcblower
    \textbf{Doğru Cevap: B} \\
    \textbf{Açıklama:} Naive Bayes, hesaplamayı kolaylaştırmak için özelliklerin birbirini etkilemediğini (bağımsız olduğunu) varsayar; bu varsayım gerçek dünyada nadiren tam olarak sağlanır.
\end{testcard}

\begin{testcard}{Soru 15}
    Gradient Boosting algoritmalarında (XGBoost, LightGBM vb.), modeller nasıl eğitilir?
    \begin{itemize}
        \item[A)] Ağaçlar birbirine paralel olarak eğitilir
        \item[B)] Her yeni ağaç, bir önceki ağacın yaptığı hataları (artıkları) azaltacak şekilde sıralı eğitilir
        \item[C)] Ağaçlar rastgele seçilir
        \item[D)] Sadece tek bir ağaç kullanılır
    \end{itemize}
    \tcblower
    \textbf{Doğru Cevap: B} \\
    \textbf{Açıklama:} Boosting, zayıf öğrenicileri (weak learners) ardışıl olarak ekleyerek her adımda toplam hatayı minimize etmeye çalışır.
\end{testcard}

\begin{testcard}{Soru 16}
    Veri seti dairesel veya karmaşık şekilli kümeleri tespit etmekte K-Means'ten daha başarılı olan, yoğunluk tabanlı kümeleme algoritması hangisidir?
    \begin{itemize}
        \item[A)] Hiyerarşik Kümeleme
        \item[B)] DBSCAN
        \item[C)] PCA
        \item[D)] Karar Ağaçları
    \end{itemize}
    \tcblower
    \textbf{Doğru Cevap: B} \\
    \textbf{Açıklama:} DBSCAN, noktaların birbirine yakınlık yoğunluğuna odaklanır; bu da onun dairesel olmayan kümeleri ve gürültüyü (noise) tespit etmesini sağlar.
\end{testcard}

\begin{testcard}{Soru 17}
    Denetimsiz öğrenmede, veriler arasındaki hiyerarşiyi bir "Dendrogram" (ağaç diyagramı) yardımıyla gösteren yöntem hangisidir?
    \begin{itemize}
        \item[A)] Lojistik Regresyon
        \item[B)] Hiyerarşik Kümeleme (Hierarchical Clustering)
        \item[C)] Rastgele Orman
        \item[D)] k-NN
    \end{itemize}
    \tcblower
    \textbf{Doğru Cevap: B} \\
    \textbf{Açıklama:} Dendrogram, verilerin hangi benzerlik seviyelerinde birleştiğini en aşağıdan yukarıya doğru bir ağaç yapısıyla gösterir.
\end{testcard}

\begin{testcard}{Soru 18}
    Boyut indirgemede kullanılan LDA (Doğrusal Diskriminant Analizi) algoritmasını PCA'dan ayıran temel özellik nedir?
    \begin{itemize}
        \item[A)] LDA'nın denetimli (sınıf etiketlerini kullanan) bir yöntem olması
        \item[B)] LDA'nın sadece regresyonda kullanılması
        \item[C)] LDA'nın daha yavaş çalışması
        \item[D)] LDA'nın gürültüyü artıran bir yöntem olması
    \end{itemize}
    \tcblower
    \textbf{Doğru Cevap: A} \\
    \textbf{Açıklama:} PCA verinin genel yapısını korurken etiketlere bakmaz; LDA ise sınıfları birbirinden ayırmak için verideki etiket bilgisini (labels) kullanır.
\end{testcard}

\begin{testcard}{Soru 19}
    Bir regresyon modelinde, gerçek değerler ile tahmin edilen değerler arasındaki farkların karelerinin ortalamasına ne ad verilir?
    \begin{itemize}
        \item[A)] Ortalama Mutlak Hata (MAE)
        \item[B)] Ortalama Kare Hata (MSE)
        \item[C)] F1 Skoru
        \item[D)] Hassasiyet (Precision)
    \end{itemize}
    \tcblower
    \textbf{Doğru Cevap: B} \\
    \textbf{Açıklama:} MSE (Mean Squared Error), hataların karesini aldığı için büyük hataları küçük hatalara göre çok daha sert cezalandırır.
\end{testcard}

\begin{testcard}{Soru 20}
    Karar ağaçlarında, aşırı öğrenmeyi (overfitting) engellemek için ağacın gereksiz dallarının temizlenmesi işlemine ne denir?
    \begin{itemize}
        \item[A)] Budama (Pruning)
        \item[B)] Dallanma (Branching)
        \item[C)] Ölçeklendirme (Scaling)
        \item[D)] Normalizasyon
    \end{itemize}
    \tcblower
    \textbf{Doğru Cevap: A} \\
    \textbf{Açıklama:} Budama, ağacı basitleştirerek modelin eğitim verisini ezberlemesini önler ve genelleme gücünü artırır.
\end{testcard}

\begin{testcard}{Soru 21}
    Yapay Sinir Ağlarında, hatanın çıktıdan girdiye doğru yayılarak ağırlıkların güncellendiği temel öğrenme mekanizması hangisidir?
    \begin{itemize}
        \item[A)] İleri Besleme (Forward Propagation)
        \item[B)] Geri Yayılım (Backpropagation)
        \item[C)] Aktivasyon
        \item[D)] Pooling
    \end{itemize}
    \tcblower
    \textbf{Doğru Cevap: B} \\
    \textbf{Açıklama:} Backpropagation, gradyan inişi ile ağırlıkların hatayı azaltacak şekilde optimize edilmesini sağlayan temel eğitim mekanizmasıdır.
\end{testcard}

\begin{testcard}{Soru 22}
    Temel Bileşen Analizi (PCA) ve Doğrusal Diskriminant Analizi (LDA) arasındaki en temel fark aşağıdakilerden hangisidir?
    \begin{itemize}
        \item[A)] LDA doğrusal bir dönüşüm yaparken, PCA yalnızca doğrusal olmayan verilerde çalışır.
        \item[B)] PCA denetimsiz bir yöntemken, LDA sınıf etiketlerini kullanan denetli bir yöntemdir.
        \item[C)] LDA bileşenler arasında korelasyon bırakırken, PCA bileşenleri birbirinden bağımsız hale getirir.
        \item[D)] PCA gürültü azaltma için kullanılırken, LDA yalnızca veri sıkıştırma için kullanılır.
    \end{itemize}
    \tcblower
    \textbf{Doğru Cevap: B} \\
    \textbf{Açıklama:} PCA verinin toplam varyansına (bilgisine) odaklanırken, LDA sınıfları birbirinden en uzağa itecek (separability) yönü bulur.
\end{testcard}

\begin{testcard}{Soru 23}
    PCA algoritmasında verinin hangi yöne projekte edileceğine karar vermek için hangi matematiksel yapıdan yararlanılır?
    \begin{itemize}
        \item[A)] Kovaryans matrisinin özvektörleri (eigenvectors)
        \item[B)] Veri matrisinin standart sapma değerleri
        \item[C)] Sigmoid fonksiyonunun türevi
        \item[D)] Sınıflar arası saçılım matrisinin izi (trace)
    \end{itemize}
    \tcblower
    \textbf{Doğru Cevap: A} \\
    \textbf{Açıklama:} Özvektörler verideki varyansın (bilginin) yönünü, özdeğerler ise bu yönün önem derecesini belirtir.
\end{testcard}

\begin{testcard}{Soru 24}
    Kümeleme analizinde kullanılan 'Dirsek (Elbow) Yöntemi' temel olarak neyi belirlemek için kullanılır?
    \begin{itemize}
        \item[A)] Hiyerarşik kümelemedeki en uzun mesafe eşiğini
        \item[B)] Optimum küme sayısı olan $K$ değerini
        \item[C)] Küme merkezlerinin başlangıç koordinatlarını
        \item[D)] Verideki aykırı değerlerin (outliers) miktarını
    \end{itemize}
    \tcblower
    \textbf{Doğru Cevap: B} \\
    \textbf{Açıklama:} Dirsek yöntemi, küme sayısı arttıkça toplam hata düşüşünün yavaşladığı "dirsek" noktasını bularak en verimli K değerini seçer.
\end{testcard}

\begin{testcard}{Soru 25}
    Regresyon modellerinde Ortalama Kare Hata (MSE) metriğinin Ortalama Mutlak Hata (MAE) metriğine göre en belirgin dezavantajı nedir?
    \begin{itemize}
        \item[A)] Aykırı değerlere (outliers) karşı aşırı duyarlı olması ve hatayı cezalandırması
        \item[B)] Matematiksel olarak türevinin alınmasının zor olması
        \item[C)] Modelin açıklanan varyansını ölçememesi
        \item[D)] Sonucun hedef değişkenle aynı birimde olmaması
    \end{itemize}
    \tcblower
    \textbf{Doğru Cevap: A} \\
    \textbf{Açıklama:} MSE hataların karesini aldığı için, büyük sapmalar skoru dramatik şekilde artırır; bu da aykırı değerlerin modeli domine etmesine yol açabilir.
\end{testcard}

\begin{testcard}{Soru 26}
    Lojistik Regresyon modelinde doğrusal kombinasyonun ($z = w^T x + b$) sonucunu 0 ile 1 arasında bir olasılık değerine dönüştüren fonksiyon hangisidir?
    \begin{itemize}
        \item[A)] ReLU Fonksiyonu: $g(z) = \max(0, z)$
        \item[B)] Sigmoid Fonksiyonu: $\sigma(z) = \frac{1}{1 + e^{-z}}$
        \item[C)] Softmax Fonksiyonu
        \item[D)] Minkowski Uzaklığı
    \end{itemize}
    \tcblower
    \textbf{Doğru Cevap: B} \\
    \textbf{Açıklama:} Sigmoid fonksiyonu, herhangi bir gerçel sayı değerini 0 ile 1 arasına sıkıştırarak sınıf olasılığı üretir.
\end{testcard}

\begin{testcard}{Soru 27}
    Yapay Sinir Ağlarında 'Geri Yayılım' (Backpropagation) algoritmasının temel amacı nedir?
    \begin{itemize}
        \item[A)] Girdi verilerini bir katmandan diğerine doğrusal olmayan şekilde aktarmak
        \item[B)] Aşırı öğrenmeyi (overfitting) önlemek için nöronları rastgele kapatmak
        \item[C)] Zincir kuralı kullanarak kaybın ağırlıklara göre gradyanını hesaplamak ve ağırlıkları güncellemek
        \item[D)] Verideki eksik özellikleri istatistiksel yöntemlerle tamamlamak
    \end{itemize}
    \tcblower
    \textbf{Doğru Cevap: C} \\
    \textbf{Açıklama:} Geri yayılım, gradyan inişi (gradient descent) ile modelin ağırlıklarını hatayı azaltacak yönde optimize etmek için kullanılır.
\end{testcard}

\begin{testcard}{Soru 28}
    Destek Vektör Makineleri'nde (SVM) 'Kernel Trick' kullanımı hangi senaryoda gereklidir?
    \begin{itemize}
        \item[A)] Sınıf dengesizliği (class imbalance) problemi yaşandığında
        \item[B)] Eğitim verisi setinde çok fazla gürültü ve eksik veri olduğunda
        \item[C)] Veri noktaları mevcut boyutta doğrusal bir hiper-düzlem ile ayrılamadığında
        \item[D)] Modelin eğitim süresini kısaltmak ve hesaplama maliyetini düşürmek istendiğinde
    \end{itemize}
    \tcblower
    \textbf{Doğru Cevap: C} \\
    \textbf{Açıklama:} Kernel hilesi, veriyi yüksek boyutlu bir uzaya izdüşürerek doğrusal olmayan verilerin ayrılabilir hale gelmesini sağlar.
\end{testcard}

\begin{testcard}{Soru 29}
    Düzeltilmiş R-Kare (Adjusted $R^2$) metriğinin standart R-Kare metriğinden üstünlüğü nedir?
    \begin{itemize}
        \item[A)] Modelin tahmin hızını matematiksel olarak ölçebilmesi
        \item[B)] Sadece lojistik regresyon gibi sınıflandırma modellerinde kullanılabilmesi
        \item[C)] Her zaman 0 ile 1 arasında bir değer alması
        \item[D)] Modele eklenen ve anlamlı katkısı olmayan 'gürültü' özellikler için modeli cezalandırması
    \end{itemize}
    \tcblower
    \textbf{Doğru Cevap: D} \\
    \textbf{Açıklama:} Adjusted R2, sadece modele gerçekten bilgi katan değişkenler eklendiğinde yükselir; gereksiz değişken eklenmesi durumunda ise cezalandırma yaparak düşer.
\end{testcard}

\begin{testcard}{Soru 30}
    DBSCAN kümeleme algoritmasının K-Means algoritmasına göre en önemli avantajı aşağıdakilerden hangisidir?
    \begin{itemize}
        \item[A)] Bayes teoremini temel alarak yumuşak (soft) kümeleme yapması
        \item[B)] Küme sayısını ($K$) kullanıcıdan parametre olarak alması
        \item[C)] Büyük veri setlerinde çok daha hızlı ve düşük bellek maliyetiyle çalışması
        \item[D)] Karmaşık ve düzensiz şekilli kümeleri tespit edebilmesi ve aykırı değerleri gürültü olarak işaretleyebilmesi
    \end{itemize}
    \tcblower
    \textbf{Doğru Cevap: D} \\
    \textbf{Açıklama:} DBSCAN, K-Means'in aksine küme sayısını önceden bilmeye gerek duymaz ve hem yoğunluk tabanlı karmaşık şekilleri bulabilir hem de aykırı değerleri (noise) başarıyla eleyebilir.
\end{testcard}

\begin{testcard}{Soru 31}
    K-Katlı Çapraz Doğrulama (k-Fold Cross Validation) yönteminin temel avantajı aşağıdakilerden hangisidir?
    \begin{itemize}
        \item[A)] Modelin eğitim süresini k kat azaltması
        \item[B)] Veri setindeki tüm örneklerin hem eğitim hem de test aşamasında kullanılarak verinin daha verimli değerlendirilmesi
        \item[C)] Sadece doğrusal modellerde en yüksek doğruluğu garanti etmesi
        \item[D)] Veri setindeki aykırı değerleri otomatik olarak temizlemesi
    \end{itemize}
    \tcblower
    \textbf{Doğru Cevap: B} \\
    \textbf{Açıklama:} Çapraz doğrulama, veriyi k parçaya bölüp her birini test seti olarak kullanarak modelin performansının tüm veri seti üzerinden daha güvenilir şekilde ölçülmesini sağlar.
\end{testcard}

\begin{testcard}{Soru 32}
    Model doğrulama sürecinde veri seti genellikle hangi oranlarda eğitim, doğrulama ve test setlerine ayrılır?
    \begin{itemize}
        \item[A)] \%80 Eğitim, \%10 Doğrulama, \%10 Test
        \item[B)] \%70 Eğitim, \%30 Test
        \item[C)] \%50 Eğitim, \%25 Doğrulama, \%25 Test
        \item[D)] \%60 Eğitim, \%20 Doğrulama, \%20 Test
    \end{itemize}
    \tcblower
    \textbf{Doğru Cevap: A} \\
    \textbf{Açıklama:} Verinin büyük kısmı öğrenme için ayrılırken, hiperparametre ayarı (validation) ve son tarafsız ölçüm (test) için küçük ama dengeli paylar bırakılır.
\end{testcard}

\begin{testcard}{Soru 33}
    Bir modelin eğitim verisindeki gürültüleri ezberleyip, doğrulama setinde başarısız olması durumu aşağıdakilerden hangisidir?
    \begin{itemize}
        \item[A)] Eksik Öğrenme (Underfitting)
        \item[B)] Tam Kararında (Optimum)
        \item[C)] Yüksek Yanlılık (High Bias)
        \item[D)] Aşırı Öğrenme (Overfitting)
    \end{itemize}
    \tcblower
    \textbf{Doğru Cevap: D} \\
    \textbf{Açıklama:} Aşırı öğrenme (Overfitting), modelin verideki örüntüler yerine rastgele gürültüleri ezberlemesi sonucu genelleme yeteneğini kaybetmesidir.
\end{testcard}

\begin{testcard}{Soru 34}
    Yanlılık ve Varyans dengesiyle ilgili aşağıdakilerden hangisi doğrudur?
    \begin{itemize}
        \item[A)] Yanlılık (Bias), basit modelden kaynaklanan sistematik hatadır.
        \item[B)] Düşük varyans, modelin verideki gürültüyü ezberlediği anlamına gelir.
        \item[C)] Doğrulama süreci sadece yanlılık hatasını sıfırlamayı amaçlar.
        \item[D)] Varyans arttıkça model daha kararlı hale gelir.
    \end{itemize}
    \tcblower
    \textbf{Doğru Cevap: A} \\
    \textbf{Açıklama:} Yanlılık (Bias), modelin verideki temel ilişkiyi kavrayamayacak kadar basit olmasından kaynaklanan hatadır.
\end{testcard}

\begin{testcard}{Soru 35}
    LOOCV (Leave-One-Out Cross Validation) yöntemi için aşağıdakilerden hangisi bir dezavantajdır?
    \begin{itemize}
        \item[A)] Veri israfı
        \item[B)] Yüksek yanlılık (Bias)
        \item[C)] Yüksek hesaplama maliyeti
        \item[D)] Rastgelelikten kaynaklanan yüksek belirsizlik
    \end{itemize}
    \tcblower
    \textbf{Doğru Cevap: C} \\
    \textbf{Açıklama:} LOOCV, veri setindeki her bir gözlem için ayrı bir model eğittiğinden, büyük veri setlerinde işlem süresi aşırı derecede artar.
\end{testcard}

\begin{testcard}{Soru 36}
    Sektör standardı olarak kabul edilen k-Katlı Çapraz Doğrulama (k-Fold CV) için yaygın olarak tercih edilen $k$ değerleri nelerdir?
    \begin{itemize}
        \item[A)] $k = n$ (Gözlem sayısı)
        \item[B)] $k = 100$
        \item[C)] $k = 2$ veya $k = 3$
        \item[D)] $k = 5$ veya $k = 10$
    \end{itemize}
    \tcblower
    \textbf{Doğru Cevap: D} \\
    \textbf{Açıklama:} k=5 veya k=10 değerleri, varyans ve hesaplama maliyeti arasındaki en dengeli sonuçları verdiği için standart olarak kabul edilir.
\end{testcard}

\begin{testcard}{Soru 37}
    İç İçe (Nested) CV yönteminin temel kullanım amacı nedir?
    \begin{itemize}
        \item[A)] Küçük veri setlerinde veri üretimini artırmak
        \item[B)] Aynı anda hem hiperparametre ayarlamak hem de tarafsız performans ölçmek
        \item[C)] Veri setindeki sınıf dengesizliğini gidermek
        \item[D)] İki farklı algoritmanın performansını istatistiksel olarak karşılaştırmak
    \end{itemize}
    \tcblower
    \textbf{Doğru Cevap: B} \\
    \textbf{Açıklama:} İç içe çapraz doğrulamada iç döngü parametre seçer, dış döngü ise seçilen modelin başarısını tarafsız bir şekilde test eder.
\end{testcard}

\begin{testcard}{Soru 38}
    Bootstrap (Önyükleme) yöntemini diğer doğrulama yöntemlerinden ayıran temel özellik nedir?
    \begin{itemize}
        \item[A)] Hata payının sadece eğitim seti üzerinden hesaplanması
        \item[B)] Sınıf oranlarının korunması
        \item[C)] Veriden yerine koyarak (with replacement) örnek çekilmesi
        \item[D)] Verinin sadece bir kez bölünmesi
    \end{itemize}
    \tcblower
    \textbf{Doğru Cevap: C} \\
    \textbf{Açıklama:} Bootstrap yönteminde orijinal veri setinden, aynı örneklerin tekrar seçilebildiği yeni alt kümeler oluşturulur.
\end{testcard}

\begin{testcard}{Soru 39}
    Sınıflandırma problemlerinde başarının ölçülmesi için kullanılan 'Confusion Matrix' yapısında, gerçekte pozitif olup model tarafından negatif tahmin edilen değer hangisidir?
    \begin{itemize}
        \item[A)] False Positive (Yanlış Pozitif)
        \item[B)] False Negative (Yanlış Negatif)
        \item[C)] True Positive (Doğru Pozitif)
        \item[D)] True Negative (Doğru Negatif)
    \end{itemize}
    \tcblower
    \textbf{Doğru Cevap: B} \\
    \textbf{Açıklama:} Modelin gerçek bir vakayı kaçırması durumudur; "Pozitif" olan veri "Negatif" olarak hatalı (False) tahmin edilmiştir.
\end{testcard}

\begin{testcard}{Soru 40}
    Regresyon problemlerinde model başarısını değerlendirmek için kullanılan temel metrik aşağıdakilerden hangisidir?
    \begin{itemize}
        \item[A)] F1 Skoru
        \item[B)] Doğruluk (Accuracy)
        \item[C)] Duyarlılık (Recall)
        \item[D)] Ortalama Kare Hata (MSE)
    \end{itemize}
    \tcblower
    \textbf{Doğru Cevap: D} \\
    \textbf{Açıklama:} Sayısal değer tahmin edilen regresyon modellerinde hata, gerçek ve tahmin edilen değerler arasındaki farkın karesi (MSE) ile ölçülür.
\end{testcard}

\begin{testcard}{Soru 41}
    Karar ağacı rehberine göre, elinizde çok küçük bir veri seti varsa hangi yöntemi seçmeniz önerilir?
    \begin{itemize}
        \item[A)] LOOCV
        \item[B)] 5x2 Çapraz Doğrulama
        \item[C)] Basit Doğrulama Seti
        \item[D)] Bootstrap
    \end{itemize}
    \tcblower
    \textbf{Doğru Cevap: A} \\
    \textbf{Açıklama:} Veri seti çok küçükse, her bir gözlemden maksimum faydayı sağlamak için her seferinde sadece bir örneği dışarıda bırakan LOOCV önerilir.
\end{testcard}

\begin{testcard}{Soru 42}
    Makine öğrenmesi modellerinde "İç Ağırlıklar" ile "Hiperparametreler" arasındaki temel fark nedir?
    \begin{itemize}
        \item[A)] Hiperparametreler eğitim sırasında veriden öğrenilir, iç ağırlıklar ise kullanıcı tarafından manuel olarak girilir.
        \item[B)] İç ağırlıklar eğitim sırasında veriden öğrenilirken, hiperparametreler dışarıdan ayarlanan kontrol değişkenleridir.
        \item[C)] İç ağırlıklar modelin karmaşıklığını belirlerken, hiperparametreler sadece tahmin hızını etkiler.
        \item[D)] Hiperparametreler sadece k-NN gibi basit algoritmalarda kullanılır, iç ağırlıklar ise tüm algoritmalarda bulunur.
    \end{itemize}
    \tcblower
    \textbf{Doğru Cevap: B} \\
    \textbf{Açıklama:} Ağırlıklar modelin veriden öğrendiği parametrelerdir; hiperparametreler ise k (komşu sayısı) veya öğrenme hızı gibi kullanıcı ayarlarıdır.
\end{testcard}

\begin{testcard}{Soru 43}
    Grid Search (Izgara Araması) yönteminin hesaplama maliyeti $O(L^F)$ formülü ile ifade edilir. Buradaki $F$ değişkeni neyi temsil etmektedir?
    \begin{itemize}
        \item[A)] Optimize edilen toplam faktör (hiperparametre) sayısı.
        \item[B)] Veri setindeki katlama (fold) sayısı.
        \item[C)] Her faktörün sahip olduğu seviye sayısı.
        \item[D)] Toplam eğitilecek model sayısı.
    \end{itemize}
    \tcblower
    \textbf{Doğru Cevap: A} \\
    \textbf{Açıklama:} F (Factor), ayarlanmak istenen hiperparametre sayısını; L (Level) ise her parametre için denenecek değer sayısını temsil eder.
\end{testcard}

\begin{testcard}{Soru 44}
    Aynı model bütçesi (örneğin 25 deneme) kullanıldığında, Random Search'ün Grid Search'e göre ana stratejik avantajı nedir?
    \begin{itemize}
        \item[A)] Deneme sayısını her zaman 25'ten daha az tutması
        \item[B)] Önemli hiperparametreler üzerinde daha fazla benzersiz değer deneme olasılığının olması
        \item[C)] Her zaman mutlak küresel optimumu bulmayı garanti etmesi
        \item[D)] Sadece tek bir hiperparametre varken daha hızlı olması
    \end{itemize}
    \tcblower
    \textbf{Doğru Cevap: B} \\
    \textbf{Açıklama:} Random Search, parametre uzayında rastgele noktalar seçtiği için, önemli parametre eksenlerinde ızgara aramasına göre daha fazla farklı değer yakalayabilir.
\end{testcard}

\begin{testcard}{Soru 45}
    Bayesyen Optimizasyon yöntemini Random Search'ten ayıran temel fark nedir?
    \begin{itemize}
        \item[A)] Daha fazla bellek kullanması
        \item[B)] Her denemeyi rastgele seçmek yerine geçmiş denemelerin sonuçlarından öğrenerek bir sonraki adımı planlaması
        \item[C)] Sadece görsel verilerde çalışabilmesi
        \item[D)] Hiçbir matematiksel model kullanmaması
    \end{itemize}
    \tcblower
    \textbf{Doğru Cevap: B} \\
    \textbf{Açıklama:} Bayesyen optimizasyon, amaç fonksiyonunun bir olasılıksal modelini (surrogate model) kurarak aramayı akıllıca yönlendirir.
\end{testcard}

\begin{testcard}{Soru 46}
    Aşağıdakilerden hangisi Grid Search yönteminin "Garantici" karakteristiğini en iyi açıklar?
    \begin{itemize}
        \item[A)] Belirlenen parametre kümesi üzerindeki tüm kombinasyonları istisnasız test etmesi.
        \item[B)] Rastgele seçimlerle en iyi noktayı tesadüfen bulması.
        \item[C)] Daima küresel zirveyi (global optimum) bulmayı vaat etmesi.
        \item[D)] Hesaplama maliyetini her zaman minimumda tutması.
    \end{itemize}
    \tcblower
    \textbf{Doğru Cevap: A} \\
    \textbf{Açıklama:} Grid Search, tanımlanan aralıktaki hiçbir noktayı atlamadan "kaba kuvvet" (brute force) ile tarama yapar.
\end{testcard}

\begin{testcard}{Soru 47}
    Genetik Algoritmalar, parametre kombinasyonlarını iyileştirmek için hangi iki temel mekanizmayı kullanır?
    \begin{itemize}
        \item[A)] Izgara tarama ve rastgele örnekleme.
        \item[B)] Isıtma ve soğutma döngüleri.
        \item[C)] Mutasyon ve çaprazlama.
        \item[D)] Doğrusal regresyon ve gradyan inişi.
    \end{itemize}
    \tcblower
    \textbf{Doğru Cevap: C} \\
    \textbf{Açıklama:} En iyi parametreleri "ebeveyn" olarak seçip mutasyon ve çaprazlama ile yeni nesillere daha başarılı kombinasyonlar aktarmaya çalışır.
\end{testcard}

\begin{testcard}{Soru 48}
    SVM (Destek Vektör Makineleri) algoritmasında kullanılan $C$ parametresi ne tür bir değişkendir?
    \begin{itemize}
        \item[A)] Bir model hiperparametresidir.
        \item[B)] Bir veri temizleme yöntemidir.
        \item[C)] Veriden öğrenilen bir iç ağırlıktır.
        \item[D)] Bir performans metriğidir.
    \end{itemize}
    \tcblower
    \textbf{Doğru Cevap: A} \\
    \textbf{Açıklama:} C parametresi, modelin marjin genişliği ile hata toleransı arasındaki dengeyi kontrol eden kullanıcı ayarıdır.
\end{testcard}

\begin{testcard}{Soru 49}
    Hiperparametre arama uzayında "Hesaplama Darboğazı" (Computational Bottleneck) terimi ne zaman ortaya çıkar?
    \begin{itemize}
        \item[A)] Veri seti çok küçük olduğunda.
        \item[B)] Grid Search yönteminde faktör veya seviye sayısı arttığında.
        \item[C)] Random Search bütçesi çok düşük tutulduğunda.
        \item[D)] Çapraz doğrulama katlamaları tek bir grupta toplandığında.
    \end{itemize}
    \tcblower
    \textbf{Doğru Cevap: B} \\
    \textbf{Açıklama:} Grid Search her yeni parametre veya seviye eklendiğinde model sayısını üstel olarak artırdığı için sistem kaynaklarını tüketir.
\end{testcard}

\newpage
\subsection{Model Testi}

\begin{testcard}{Soru 50}
    Bir makine öğrenmesi modelinin toplam hata miktarını ($MSE$) belirleyen bileşenler düşünüldüğünde, model iyileştirilerek azaltılamayan hata türü aşağıdakilerden hangisidir?
    \begin{itemize}
        \item[A)] Varyans (Variance)
        \item[B)] Eğitim Hatası (Training Error)
        \item[C)] İndirgenemez Hata ($\sigma^2$ - Gürültü)
        \item[D)] Yanlılık (Bias)
    \end{itemize}
    \tcblower
    \textbf{Doğru Cevap: C} \\
    \textbf{Açıklama:} İndirgenemez hata, verinin doğasından gelen gürültüdür; model ne kadar karmaşık olursa olsun bu hata payı sıfırlanamaz.
\end{testcard}

\begin{testcard}{Soru 51}
    Izgara Araması (Grid Search) stratejisinde, her bir hiperparametrenin seviye sayısı $L$ ve toplam hiperparametre sayısı $F$ ise, hesaplama maliyeti hangi formülle ifade edilir?
    \begin{itemize}
        \item[A)] $O(L \cdot F)$
        \item[B)] $O(\log(L^F))$
        \item[C)] $O(L^F)$
        \item[D)] $O(F^L)$
    \end{itemize}
    \tcblower
    \textbf{Doğru Cevap: C} \\
    \textbf{Açıklama:} Grid Search'te her parametrenin her seviyesi birbiriyle eşleştiği için toplam kombinasyon sayısı üsteldir.
\end{testcard}

\begin{testcard}{Soru 52}
    Bir ölçümün matematiksel olarak 'metrik' sayılabilmesi için üç temel kuralı sağlaması gerekir. Aşağıdakilerden hangisi bu kurallardan biridir?
    \begin{itemize}
        \item[A)] Normalizasyon ([0, 1] aralığı)
        \item[B)] Üçgen Eşitsizliği ($d(x, y) \leq d(x, z) + d(z, y)$)
        \item[C)] Yön Bağımlılığı
        \item[D)] Kosinüs Benzerliği
    \end{itemize}
    \tcblower
    \textbf{Doğru Cevap: B} \\
    \textbf{Açıklama:} Bir metrik; negatif olmama, simetri ve üçgen eşitsizliği (doğrudan yol dolaylı yoldan uzun olamaz) kurallarını sağlamalıdır.
\end{testcard}

\begin{testcard}{Soru 53}
    Yetersiz Uyum (Underfitting) yaşayan bir modeli iyileştirmek için aşağıdakilerden hangisi uygun bir çözüm stratejisidir?
    \begin{itemize}
        \item[A)] Erken Durdurma (Early Stopping) uygulamak
        \item[B)] Düzenlileştirme (Regularization) katsayısını ($\lambda$) artırmak
        \item[C)] Daha fazla eğitim verisi toplamak
        \item[D)] Modelin esnekliğini/karmaşıklığını artırmak
    \end{itemize}
    \tcblower
    \textbf{Doğru Cevap: D} \\
    \textbf{Açıklama:} Underfitting, modelin veriyi kavrayamayacak kadar basit olmasından kaynaklanır; çözüm modelin karmaşıklığını artırmaktır.
\end{testcard}

\begin{testcard}{Soru 54}
    Kosinüs Benzerliği (Cosine Similarity) ölçütü kullanılırken odak noktası aşağıdakilerden hangisidir?
    \begin{itemize}
        \item[A)] Vektörlerin bileşenlerinin toplamı
        \item[B)] İki nokta arasındaki en kısa kuş uçuşu mesafe
        \item[C)] Vektörlerin mutlak büyüklükleri (Magnitude)
        \item[D)] Vektörlerin yönü ve aralarındaki açı
    \end{itemize}
    \tcblower
    \textbf{Doğru Cevap: D} \\
    \textbf{Açıklama:} Kosinüs benzerliği vektörlerin uzunluğundan bağımsızdır, sadece aralarındaki açıya bakarak yönsel benzerliği ölçer.
\end{testcard}

\begin{testcard}{Soru 55}
    Rastgele Arama (Random Search) stratejisinin, Izgara Araması'na (Grid Search) göre yüksek boyutlu arama uzaylarında daha avantajlı olmasının temel sebebi nedir?
    \begin{itemize}
        \item[A)] Hesaplama maliyetinin üstel olarak artması
        \item[B)] Önemli hiperparametreler üzerinde daha fazla benzersiz değer denemesi
        \item[C)] Her zaman mutlak küresel optimumu garanti etmesi
        \item[D)] Sadece doğrusal modellerde çalışması
    \end{itemize}
    \tcblower
    \textbf{Doğru Cevap: B} \\
    \textbf{Açıklama:} Random Search, önemsiz parametrelere vakit harcamak yerine, rastgele seçimlerle önemli parametre eksenlerini daha yoğun tarar.
\end{testcard}

\begin{testcard}{Soru 56}
    Aşırı Uyum (Overfitting) problemi yaşayan bir modelde Yanlılık (Bias) ve Varyans (Variance) değerlerinin durumu nasıldır?
    \begin{itemize}
        \item[A)] Hem Yanlılık hem Varyans çok yüksektir
        \item[B)] Düşük Yanlılık, Yüksek Varyans
        \item[C)] Yüksek Yanlılık, Düşük Varyans
        \item[D)] Hem Yanlılık hem Varyans çok düşüktür
    \end{itemize}
    \tcblower
    \textbf{Doğru Cevap: B} \\
    \textbf{Açıklama:} Aşırı uyumda model veriye çok iyi oturur (Düşük Yanlılık) ancak yeni verilere karşı aşırı tepki verir (Yüksek Varyans).
\end{testcard}

\begin{testcard}{Soru 57}
    Manhattan Mesafesi ($L_1$ - City Block), hangi veri özelliklerine sahip durumlarda Öklid Mesafesine ($L_2$) göre daha dirençli ve uygundur?
    \begin{itemize}
        \item[A)] Yüksek boyutlu uzaylarda ve ayrık (discrete) özelliklerde
        \item[B)] Sürekli ve düşük boyutlu verilerde
        \item[C)] Sadece metin verilerinin yönsel karşılaştırmasında
        \item[D)] Veri setinde hiç aykırı değer bulunmadığında
    \end{itemize}
    \tcblower
    \textbf{Doğru Cevap: A} \\
    \textbf{Açıklama:} Yüksek boyutlu veya seyrek (sparse) verilerde Manhattan mesafesi, Öklid mesafesine göre daha kararlı sonuçlar verir.
\end{testcard}

\begin{testcard}{Soru 58}
    Regresyon modellerinde kullanılan $AIC$ (Akaike Bilgi Kriteri) ve $BIC$ (Bayesian Bilgi Kriteri) için aşağıdakilerden hangisi doğrudur?
    \begin{itemize}
        \item[A)] Daha küçük değer, daha az karmaşıklıkla daha iyi uyum sağlayan modeli gösterir.
        \item[B)] Sadece eğitim hatasına odaklanırlar, model karmaşıklığını görmezden gelirler.
        \item[C)] Değer ne kadar büyükse model o kadar başarılıdır.
        \item[D)] AIC ve BIC değerleri her zaman negatif olmak zorundadır.
    \end{itemize}
    \tcblower
    \textbf{Doğru Cevap: A} \\
    \textbf{Açıklama:} Bu kriterler model başarısını karmaşıklıkla cezalandırır; dolayısıyla daha düşük skor, daha "optimal" (yeterince basit ve başarılı) modeli gösterir.
\end{testcard}

\begin{testcard}{Soru 59}
    Jaccard Benzerliği, özellikle hangi tür veri yapılarının karşılaştırılmasında merkezde yer alır?
    \begin{itemize}
        \item[A)] Kümelerin (sets) örtüşme miktarı
        \item[B)] Vektörlerin uzaydaki açısal farkı
        \item[C)] Zaman serilerindeki trend yönü
        \item[D)] Sayısal vektörlerin büyüklük farkı
    \end{itemize}
    \tcblower
    \textbf{Doğru Cevap: A} \\
    \textbf{Açıklama:} Jaccard benzerliği, iki kümenin kesişiminin birleşimine oranıdır; örneğin metin benzerliğinde kelime örtüşmelerini ölçer.
\end{testcard}

\newpage
\subsection{Regresyon Testi}

\begin{testcard}{Soru 60}
    Regresyon modellerinde 'Hata' veya 'Artık Değer' (Residual) kavramı matematiksel olarak nasıl ifade edilir?
    \begin{itemize}
        \item[A)] $Hata = y_i^2 - \hat{y}_i^2$
        \item[B)] $Hata = |y_i + \hat{y}_i|$
        \item[C)] $Hata = y_i - \hat{y}_i$
        \item[D)] $Hata = \frac{y_i}{\hat{y}_i}$
    \end{itemize}
    \tcblower
    \textbf{Doğru Cevap: C} \\
    \textbf{Açıklama:} Hata, verideki gerçek değer ile modelin yaptığı tahmin arasındaki basit farktır.
\end{testcard}

\begin{testcard}{Soru 61}
    Aykırı değerlerin (outliers) yoğun olduğu bir veri setinde model performansını ölçmek için hangi metrik 'daha stabil' bir sonuç verir?
    \begin{itemize}
        \item[A)] Ortalama Kare Hata (MSE)
        \item[B)] R-Kare ($R^2$)
        \item[C)] Kök Ortalama Kare Hata (RMSE)
        \item[D)] Ortalama Mutlak Hata (MAE)
    \end{itemize}
    \tcblower
    \textbf{Doğru Cevap: D} \\
    \textbf{Açıklama:} MAE hataların mutlak değerini aldığı için, büyük sapmaları karesel olarak cezalandırmaz ve daha dengeli bir skor verir.
\end{testcard}

\begin{testcard}{Soru 62}
    Ortalama Kare Hata (MSE) metriğinin makine öğrenmesi algoritmalarının optimizasyon aşamasında tercih edilmesinin temel nedeni nedir?
    \begin{itemize}
        \item[A)] Her zaman 0 ile 1 arasında sonuç vermesi
        \item[B)] Sonuçların hedef değişkenle aynı birimde olması
        \item[C)] Türevinin kolay alınabilmesi (differentiability)
        \item[D)] Aykırı değerleri tamamen yok sayması
    \end{itemize}
    \tcblower
    \textbf{Doğru Cevap: C} \\
    \textbf{Açıklama:} MSE, matematiksel olarak türetilebilir bir fonksiyon olduğu için gradyan inişi (gradient descent) optimizasyonunda yaygın olarak tercih edilir.
\end{testcard}

\begin{testcard}{Soru 63}
    Ridge ve LASSO regresyon modellerinde kullanılan ceza (penalty) teriminin temel amacı nedir?
    \begin{itemize}
        \item[A)] Modelin eğitim süresini kısaltmak
        \item[B)] Katsayıları (ağırlıkları) sınırlayarak aşırı öğrenmeyi (overfitting) engellemek
        \item[C)] Veri setindeki eksik değerleri tahmin etmek
        \item[D)] Sınıflandırma doğruluğunu artırmak
    \end{itemize}
    \tcblower
    \textbf{Doğru Cevap: B} \\
    \textbf{Açıklama:} Düzenlileştirme (regularization), katsayıların aşırı büyümesini engelleyerek modelin daha basit ve genellenebilir kalmasını sağlar.
\end{testcard}

\begin{testcard}{Soru 64}
    Lojistik Regresyon modelinde 'Log-Odds' (Lojit) ifadesinin doğrusal bir formülle ($w^T x$) modellenmesinin asıl amacı nedir?
    \begin{itemize}
        \item[A)] Çıktıyı doğrudan bir sınıfa atamak
        \item[B)] Olasılık değerini (0,1) tüm reel sayılar uzayına ($-\infty, +\infty$) haritalamak
        \item[C)] Modeli daha karmaşık hale getirmek
        \item[D)] Sadece kategorik değişkenlerle çalışabilmek
    \end{itemize}
    \tcblower
    \textbf{Doğru Cevap: B} \\
    \textbf{Açıklama:} Odds oranının logaritması alınarak, sınırlı olan olasılık uzayı doğrusal modellerin çalışabileceği sınırsız bir uzaya dönüştürülür.
\end{testcard}

\begin{testcard}{Soru 65}
    Dengesiz bir veri setinde (örn: \%99 Negatif, \%1 Pozitif), model başarısını ölçmek için hangisi en az güvenilir metriktir?
    \begin{itemize}
        \item[A)] Precision (Kesinlik)
        \item[B)] Recall (Duyarlılık)
        \item[C)] F1-Skoru
        \item[D)] Accuracy (Doğruluk)
    \end{itemize}
    \tcblower
    \textbf{Doğru Cevap: D} \\
    \textbf{Açıklama:} Doğruluk (Accuracy), azınlık sınıfını tamamen görmezden gelen bir modeli bile başarılı gösterebileceği için dengesiz verilerde yanıltıcıdır.
\end{testcard}

\begin{testcard}{Soru 66}
    $R^2 = 0.82$ olan bir model için aşağıdakilerden hangisi en doğru yorumdur?
    \begin{itemize}
        \item[A)] Modeldeki özelliklerin \%82'si tahmin yapmak için gereksizdir.
        \item[B)] Hedef değişkendeki değişkenliğin \%82'si modeldeki özellikler tarafından açıklanabiliyor.
        \item[C)] Model yapılan tahminlerde ortalama \%82 oranında hata yapmaktadır.
        \item[D)] Modelin tahminleri ile gerçek değerler arasındaki fark 0.82 birimdir.
    \end{itemize}
    \tcblower
    \textbf{Doğru Cevap: B} \\
    \textbf{Açıklama:} R-Kare, modelin verideki değişkenliği (varyansı) açıklama yüzdesini gösterir.
\end{testcard}

\begin{testcard}{Soru 67}
    Modele tamamen alakasız (gürültü) değişkenler eklendiğinde $R^2$ ve Düzeltilmiş $R^2$ (Adjusted $R^2$) değerleri nasıl değişir?
    \begin{itemize}
        \item[A)] Her iki değer de her zaman azalır.
        \item[B)] Alakasız değişkenler hiçbir metriği etkilemez.
        \item[C)] $R^2$ artar veya aynı kalır, Düzeltilmiş $R^2$ ise puanı düşürerek cezalandırır.
        \item[D)] $R^2$ azalır, Düzeltilmiş $R^2$ ise her zaman artmaya devam eder.
    \end{itemize}
    \tcblower
    \textbf{Doğru Cevap: C} \\
    \textbf{Açıklama:} R-Kare modele eklenen her şeyi kabul edip yükselir; Düzeltilmiş R-Kare ise değişken sayısına göre cezalandırma yapar.
\end{testcard}

\begin{testcard}{Soru 68}
    Bir model karnesinde RMSE değerinin (53.85), MAE değerinden (42.79) belirgin şekilde yüksek olması neyin kanıtıdır?
    \begin{itemize}
        \item[A)] Veri setindeki tüm hataların birbirine eşit büyüklükte olduğunun.
        \item[B)] Modelin aşırı öğrenme (overfitting) yaptığının.
        \item[C)] Modelin mükemmel bir tahmin kapasitesine sahip olduğunun.
        \item[D)] Veri setinde bazı tahminlerin çok büyük hatalara (aykırı değerlere) sahip olduğunun.
    \end{itemize}
    \tcblower
    \textbf{Doğru Cevap: D} \\
    \textbf{Açıklama:} RMSE büyük hataları kare alarak büyüteceği için, MAE ile arasındaki farkın açılması veri setindeki büyük sapmaları işaret eder.
\end{testcard}

\begin{testcard}{Soru 69}
    Scikit-learn kütüphanesini kullanarak bir regresyon modelinin başarısını (açıklanan varyans) ölçmek için hangi fonksiyon çağrılmalıdır?
    \begin{itemize}
        \item[A)] `regression_accuracy(y_true, y_pred)`
        \item[B)] `mae_score(y_true, y_pred)`
        \item[C)] `r2_score(y_true, y_pred)`
        \item[D)] `mean_squared_error(y_true, y_pred, squared=False)`
    \end{itemize}
    \tcblower
    \textbf{Doğru Cevap: C} \\
    \textbf{Açıklama:} Sklearn'de varyans açıklama skoru `r2_score` fonksiyonu ile hesaplanır.
\end{testcard}

\newpage
\subsection{Sınıflandırma Testi}

\begin{testcard}{Soru 70}
    Sınıflandırma biliminin temel amacı aşağıdakilerden hangisidir?
    \begin{itemize}
        \item[A)] Veri setindeki gürültüyü tamamen yok ederek hatasız modeller üretmek.
        \item[B)] Ham verileri analiz ederek bunları kesin kategorik kararlara dönüştürmek.
        \item[C)] Sürekli sayısal değerleri tahmin etmek için regresyon modelleri oluşturmak.
        \item[D)] Sadece görsel verileri işleyerek nesne tanıma gerçekleştirmek.
    \end{itemize}
    \tcblower
    \textbf{Doğru Cevap: B} \\
    \textbf{Açıklama:} Sınıflandırma, girdileri önceden belirlenmiş sınıflara atama işlemidir.
\end{testcard}

\begin{testcard}{Soru 71}
    Temel Hata Oranı ($E$) hesaplanırken kullanılan $E = \frac{1}{n} |\{i | \hat{v}_i \neq v_i\}|$ formülündeki $\hat{v}_i \neq v_i$ ifadesi neyi ifade eder?
    \begin{itemize}
        \item[A)] Modelin yaptığı doğru tahminlerin toplam sayısını.
        \item[B)] Modelin tahmin yaparken kullandığı güven aralığının genişliğini.
        \item[C)] Tahmin edilen değerin gerçek değerden farklı olduğu durumları.
        \item[D)] Veri setindeki toplam örnek sayısının tahmin edilen sınıflara oranını.
    \end{itemize}
    \tcblower
    \textbf{Doğru Cevap: C} \\
    \textbf{Açıklama:} Bu ifade "eşitsizlik" belirtir; yani tahmin edilen sınıfın gerçek sınıf ile aynı olmadığı durumları sayar.
\end{testcard}

\begin{testcard}{Soru 72}
    Bir spam filtresinin önemli bir iş e-postasını yanlışlıkla 'Spam' klasörüne taşıması hangi hata türüne örnektir?
    \begin{itemize}
        \item[A)] Tip I Hata (Yalancı Alarm)
        \item[B)] Doğru Negatif (True Negative)
        \item[C)] Bayes Hatası
        \item[D)] Tip II Hata (Kaçırılan Tehlike)
    \end{itemize}
    \tcblower
    \textbf{Doğru Cevap: A} \\
    \textbf{Açıklama:} Tip I Hata (False Positive), negatif olması gereken bir durumun "pozitif" olarak yanlış alarm vermesidir.
\end{testcard}

\begin{testcard}{Soru 73}
    Karmaşıklık Matrisi (Confusion Matrix) içerisinde yer alan 'Köşegen Dışı' ($Off-Diagonal$) elemanlar bize hangi bilgiyi verir?
    \begin{itemize}
        \item[A)] İki farklı hata tipinin (FP ve FN) sayılarını.
        \item[B)] Veri setindeki toplam örnek sayısını ($n$).
        \item[C)] Modelin doğru tahmin ettiği sınıfların dağılımını.
        \item[D)] Modelin tahmin yaparken duyduğu olasılıksal güven değerini.
    \end{itemize}
    \tcblower
    \textbf{Doğru Cevap: A} \\
    \textbf{Açıklama:} Matrisin ana köşegeni doğru tahminleri, köşegen dışındaki hücreler ise yanlış tahminleri gösterir.
\end{testcard}

\begin{testcard}{Soru 74}
    Hangi durumda 'Kesinlik' ($Precision$) metriğinin yüksek olması, 'Duyarlılık' ($Recall$) metriğinden daha kritik kabul edilir?
    \begin{itemize}
        \item[A)] Sınıfların veri setinde tamamen dengeli dağıldığı durumlarda.
        \item[B)] Kanser teşhisi gibi bir vakanın atlanmasının ölümcül sonuçlar doğurabileceği durumlarda.
        \item[C)] Yalancı Alarların (Yanlış Pozitif) maliyetinin çok yüksek olduğu durumlarda.
        \item[D)] Modelin tüm sınıfları birbirinden ayırt etme gücünü tek bir skorla görmek istediğimizde.
    \end{itemize}
    \tcblower
    \textbf{Doğru Cevap: C} \\
    \textbf{Açıklama:} Yanlış alarm vermenin maliyeti çok fazlaysa (örn: masum birini suçlu bulmak), kesinlik (precision) önceliklendirilir.
\end{testcard}

\begin{testcard}{Soru 75}
    1000 kişilik bir veri setinde 990 sağlıklı, 10 hasta birey bulunmaktadır. Herkese 'Sağlıklı' diyen bir modelin başarısı hakkında hangisi söylenebilir?
    \begin{itemize}
        \item[A)] Modelin F1-Skoru \%100 seviyesindedir.
        \item[B)] Model \%99 Doğruluk (Accuracy) oranına sahip olduğu için mükemmeldir.
        \item[C)] Modelin Duyarlılık (Recall) değeri \%0 olduğu için hastalık tespiti görevinde başarısızdır.
        \item[D)] Model yüksek bir Logaritmik Kayıp (Log Loss) puanı alarak ödüllendirilir.
    \end{itemize}
    \tcblower
    \textbf{Doğru Cevap: C} \\
    \textbf{Açıklama:} Model çoğunluk sınıfını tahmin ederek \%99 doğruluk alsa da, asıl amaç olan 10 hastanın hiçbirini bulamadığı için başarısızdır.
\end{testcard}

\newpage
\subsection{Performans ve Özellik Analizi}

\begin{testcard}{Soru 76}
    F1-Skoru hesaplanırken aritmetik ortalama yerine neden harmonik ortalama tercih edilir?
    \begin{itemize}
        \item[A)] Sadece Boolean sınıflandırıcılarda kullanılabildiği için.
        \item[B)] Uç değerleri (aşırı düşük puanları) şiddetli bir şekilde cezalandırmak için.
        \item[C)] Hesaplanması daha kolay ve işlem maliyeti daha düşük olduğu için.
        \item[D)] Veri setindeki toplam örnek sayısı çok fazla olduğunda daha stabil sonuç verdiği için.
    \end{itemize}
    \tcblower
    \textbf{Doğru Cevap: B} \\
    \textbf{Açıklama:} Harmonik ortalama, precision veya recall'dan biri çok düşük olduğunda skoru aşağı çeker; dengeyi zorunlu kılar.
\end{testcard}

\begin{testcard}{Soru 77}
    AUC (Area Under the Curve) metriği için aşağıdakilerden hangisi doğrudur?
    \begin{itemize}
        \item[A)] AUC değeri arttıkça modelin hata oranı ($E$) da artar.
        \item[B)] Sadece belirli bir eşik (threshold) değeri için modelin performansını ölçer.
        \item[C)] 0.5 değeri, modelin rastgele tahmin yapmaktan daha iyi bir iş çıkarmadığını gösterir.
        \item[D)] Değer 0 ile 1 arasında değil, $-\infty$ ile $+\infty$ arasındadır.
    \end{itemize}
    \tcblower
    \textbf{Doğru Cevap: C} \\
    \textbf{Açıklama:} AUC 0.5 ise model şans eseri tahmin yapıyor demektir; 1.0 değeri ise mükemmel ayrımı ifade eder.
\end{testcard}

\begin{testcard}{Soru 78}
    Neyman-Pearson metriğinde ($E_{NP} = \lambda \frac{C_{fn}}{n} + \frac{C_{fp}}{n}$), $\lambda > 1$ olarak seçilirse bu model hakkında ne söylenebilir?
    \begin{itemize}
        \item[A)] Yanlış Negatiflerin (vaka kaçırmanın) maliyeti, yanlış alarmlardan daha kritik kabul edilmektedir.
        \item[B)] Modelin toplam hata oranı her zaman matematiksel olarak minimize edilir.
        \item[C)] Model Yanlış Pozitifleri (False Positive) önlemeye daha çok odaklanır.
        \item[D)] Tüm hata türleri model için eşit derecede önemsiz hale gelir.
    \end{itemize}
    \tcblower
    \textbf{Doğru Cevap: A} \\
    \textbf{Açıklama:} Lambda katsayısı FN (vaka kaçırma) hatasını daha büyük bir sayıyla çarparak modeli bu hatayı yapmamaya zorlar.
\end{testcard}

\begin{testcard}{Soru 79}
    $K$ adet sınıfa sahip çok sınıflı bir sistemde, sadece bir doğru tahmin senaryosu varken tam olarak kaç farklı hatalı senaryo mevcuttur?
    \begin{itemize}
        \item[A)] $K(K-1)$
        \item[B)] $K^2$
        \item[C)] $K-1$
        \item[D)] $2^K$
    \end{itemize}
    \tcblower
    \textbf{Doğru Cevap: C} \\
    \textbf{Açıklama:} Toplam K sınıf varsa, bir tanesi doğru sınıf, geri kalan K-1 tanesi ise yanlış sınıf tahminleridir.
\end{testcard}

\newpage
\subsection{Özellik Testi}

\begin{testcard}{Soru 80}
    Boyutluluk Laneti (Curse of Dimensionality) model eğitimi için gereken örnek sayısı ile özellik sayısı ($d$) arasındaki ilişkiyi nasıl tanımlar?
    \begin{itemize}
        \item[A)] Özellik sayısı arttıkça gereken örnek sayısı azalır çünkü veri daha detaylıdır.
        \item[B)] Özellik sayısı ile örnek sayısı arasında matematiksel bir ilişki yoktur.
        \item[C)] Örnek sayısı özellik sayısına göre doğrusal bir artış gösterir: $N \propto d$.
        \item[D)] Gereken örnek sayısı özellik sayısıyla birlikte üstel olarak artar: $N \propto e^d$.
    \end{itemize}
    \tcblower
    \textbf{Doğru Cevap: D} \\
    \textbf{Açıklama:} Her yeni özellik uzayı genişlettiği için, veriyi temsil edecek örnek sayısının üstel bir hızla artması gerekir.
\end{testcard}

\begin{testcard}{Soru 81}
    Filtre (Filter) yöntemlerinin sarma (Wrapper) yöntemlerine göre en belirgin avantajı aşağıdakilerden hangisidir?
    \begin{itemize}
        \item[A)] Özellikler arasındaki karmaşık etkileşimleri yakalayabilmeleri.
        \item[B)] Hesaplama maliyetinin düşük olması ve hızlı çalışmaları.
        \item[C)] En yüksek tahmin doğruluğunu garanti etmeleri.
        \item[D)] Model performansını doğrudan optimize etmeleri.
    \end{itemize}
    \tcblower
    \textbf{Doğru Cevap: B} \\
    \textbf{Açıklama:} Filtre yöntemleri istatistiksel testlere dayandığı için model eğitmeden hızlıca karar verebilir; Wrapper'lar ise her seferinde model eğitir.
\end{testcard}

\begin{testcard}{Soru 82}
    Pearson Korelasyonu ($r$) kullanılarak yapılan özellik seçiminde, hangi durum özelliğin hedef değişkenle yüksek derecede ilgili olduğunu gösterir?
    \begin{itemize}
        \item[A)] $r$ değerinin negatif olması.
        \item[B)] $r$ değerinin sonsuza yaklaşması.
        \item[C)] $|r|$ değerinin 1'e yakın olması.
        \item[D)] $r$ değerinin 0 olması.
    \end{itemize}
    \tcblower
    \textbf{Doğru Cevap: C} \\
    \textbf{Açıklama:} Mutlak değeri 1'e yakın olan bir katsayı, değişkenler arasında çok güçlü bir doğrusal ilişki (pozitif/negatif) olduğunu gösterir.
\end{testcard}

\begin{testcard}{Soru 83}
    ANOVA F-Testi yönteminde, bir özelliğin sınıfları ayırt etme gücünün yüksek olduğu nasıl anlaşılır?
    \begin{itemize}
        \item[A)] Yüksek bir $F$ değeri elde edilmesiyle.
        \item[B)] Sınıf içi varyansın ($MSW$) sınıflar arası varyanstan ($MSB$) büyük olmasıyla.
        \item[C)] Özellik değerlerinin standart sapmasının 0 olmasıyla.
        \item[D)] $F$ skorunun düşük olmasıyla.
    \end{itemize}
    \tcblower
    \textbf{Doğru Cevap: A} \\
    \textbf{Açıklama:} Yüksek F skoru, sınıflar arasındaki farkın, sınıf içi değişkenlikten çok daha büyük olduğunu kanıtlar.
\end{testcard}

\begin{testcard}{Soru 84}
    Sarma (Wrapper) yöntemlerinden biri olan İleriye Doğru Seçim (Forward Selection) süreci nasıl başlar?
    \begin{itemize}
        \item[A)] Sadece hedef değişkenle korelasyonu en düşük olan özelliklerle başlar.
        \item[B)] Tüm özelliklerin bulunduğu tam küme ile başlar.
        \item[C)] Boş bir özellik kümesi ($S = \emptyset$) ile başlar.
        \item[D)] Rastgele seçilmiş bir özellik alt kümesi ile başlar.
    \end{itemize}
    \tcblower
    \textbf{Doğru Cevap: C} \\
    \textbf{Açıklama:} İleriye doğru seçimde başlangıçta hiçbir özellik yoktur; model başarısını en çok artıran özellikler tek tek kümeye eklenir.
\end{testcard}

\begin{testcard}{Soru 85}
    LASSO ($L1$ Düzenlileştirmesi) yöntemini Ridge ($L2$) yönteminden ayıran ve özellik seçimi yapmasını sağlayan temel mekanizma nedir?
    \begin{itemize}
        \item[A)] Hesaplama karmaşıklığının $O(d^2)$ olması.
        \item[B)] Sadece doğrusal olmayan modellerde kullanılabilmesi.
        \item[C)] Bazı katsayıları tam olarak sıfıra eşitleyerek değişkenleri elemesi.
        \item[D)] Katsayıları sadece küçültmesi ama asla sıfır yapmaması.
    \end{itemize}
    \tcblower
    \textbf{Doğru Cevap: C} \\
    \textbf{Açıklama:} LASSO, mutlak değer cezası kullanarak gereksiz gördüğü özelliklerin katsayılarını sıfıra indirir.
\end{testcard}

\begin{testcard}{Soru 86}
    Özellik seçiminde Mutual Information (Karşılıklı Bilgi) yönteminin Pearson korelasyonuna göre en büyük avantajı nedir?
    \begin{itemize}
        \item[A)] Sadece kategorik verilerle çalışabilmesi.
        \item[B)] Hesaplamasının çok daha hızlı olması.
        \item[C)] Sarma (Wrapper) yöntemleri sınıfına girmesi.
        \item[D)] Değişkenler arasındaki doğrusal olmayan (non-linear) ilişkileri de yakalayabilmesi.
    \end{itemize}
    \tcblower
    \textbf{Doğru Cevap: D} \\
    \textbf{Açıklama:} Korelasyon sadece doğrusal bağları görürken, Mutual Information herhangi bir istatistiksel bağımlılığı tespit edebilir.
\end{testcard}

\begin{testcard}{Soru 87}
    Öz yinelemeli Özellik Eleme ($RFE$) algoritması, özelliklerin önemini belirlemek için neyi kullanır?
    \begin{itemize}
        \item[A)] Özellikler arasındaki Pearson korelasyon matrisini.
        \item[B)] Özelliklerin tek başına varyans değerlerini.
        \item[C)] Eğitilen modelden elde edilen katsayı ağırlıkları veya önem skorlarını ($|w_j|$).
        \item[D)] Veri setindeki eksik değer oranlarını.
    \end{itemize}
    \tcblower
    \textbf{Doğru Cevap: C} \\
    \textbf{Açıklama:} RFE, modeli tüm özelliklerle eğitir ve her adımda en düşük ağırlığa (öneme) sahip özelliği kümeden atarak ilerler.
\end{testcard}

\begin{testcard}{Soru 88}
    Ağaç tabanlı modellerde Permütasyon Önemi ($MDA$), Safsızlık Azalmasına ($MDI$) göre neden daha güvenilir kabul edilir?
    \begin{itemize}
        \item[A)] Yüksek kardinaliteli (çok sayıda eşsiz değer içeren) değişkenlere karşı daha az yanlı olması.
        \item[B)] Hesaplamasının çok daha hızlı olması nedeniyle.
        \item[C)] Matematiksel olarak daha basit bir formüle sahip olması.
        \item[D)] Sadece eğitim verisi üzerinde hesaplandığı için.
    \end{itemize}
    \tcblower
    \textbf{Doğru Cevap: A} \\
    \textbf{Açıklama:} MDI, çok sayıda eşsiz sayısal değere sahip değişkenleri haksız yere önemli gösterme eğilimindedir; MDA bu yanlılığı aşar.
\end{testcard}

\begin{testcard}{Soru 89}
    Elastic Net yönteminde kullanılan $\alpha = \frac{\lambda_1}{\lambda_1 + \lambda_2}$ parametresi neyi kontrol eder?
    \begin{itemize}
        \item[A)] Veri setindeki aykırı değerlerin etkisini.
        \item[B)] Karar ağacının maksimum derinliğini.
        \item[C)] Modelin öğrenme hızını (learning rate).
        \item[D)] $L1$ (LASSO) ve $L2$ (Ridge) düzenlileştirmeleri arasındaki karışım oranını.
    \end{itemize}
    \tcblower
    \textbf{Doğru Cevap: D} \\
    \textbf{Açıklama:} Alpha parametresi, Elastic Net'in ne kadar LASSO'ya ne kadar Ridge'e benzeyeceğini belirleyen karışım oranıdır.
\end{testcard}

\begin{testcard}{Soru 90}
    Makine öğrenmesinde "Özellik Mühendisliği" (Feature Engineering) sürecinin temel amacı nedir?
    \begin{itemize}
        \item[A)] Verileri sadece görselleştirmek.
        \item[B)] Mevcut verilerden modele ek bilgi sağlayacak yeni ve daha anlamlı değişkenler üretmek.
        \item[C)] Veri setini tamamen silip yeniden oluşturmak.
        \item[D)] Modelin eğitim süresini sonsuza kadar uzatmak.
    \end{itemize}
    \tcblower
    \textbf{Doğru Cevap: B} \\
    \textbf{Açıklama:} Verideki ham bilgiyi, modelin daha kolay öğrenebileceği ve performansı artıracak yeni özelliklere dönüştürme sürecidir.
\end{testcard}

\newpage

% ==============================================================================
% BÖLÜM 4: KAYNAKÇA
% ==============================================================================
\end{document}
